\section{A. The retained family and its intrinsic classification}

Throughout this section $\mathcal P=\{c:\operatorname{Im}c<-M\}$, and $X_c$ means the actual glued space with $v_1=\varepsilon=(1,2,-4,0)^t$, $v_2=-\varepsilon'=(-1,-3,3,0)^t$ and the untwisted cusp collar. We write $d=\widehat\delta$, $e=(0,1)^t$, $E=d\gamma$, and use column vectors. Our peripheral generators are clockwise: $s_j(g_jz)=e^{-2\pi i/m_j}s_j(z)$ and $g_1g_2g_0=1$. Thus $s(g_0z)=s(z)-1$ at the cusp. All period comparisons use
\[
 A_1=\begin{pmatrix}1&0&0&0\\6&0&1&0\\-6&-1&-1&0\\-2&1&0&1\end{pmatrix},\quad
 A_2=\begin{pmatrix}1&0&0&0\\0&0&-1&0\\-6&1&0&0\\3&0&1&1\end{pmatrix},\qquad
 R_g(z)\Pi_c(z)=\Pi_c(gz)A_g.
\]
Here $R_1=\left(\begin{smallmatrix}-1/\tau&0\\(1-\mu)/\tau&1\end{smallmatrix}\right)$, $R_2=\left(\begin{smallmatrix}1/\tau&0\\-\mu/\tau&1\end{smallmatrix}\right)$, and $R_0=I$. The distinction between source and target markings in each equation below is retained by its displayed period equation.

\subsection{A.1. Canonical sections and the finite covers}

\begin{proposition}[A.1.1; Proved]\label{st:A.1.1}
For the constructed holomorphic map $f_c:X_c\to\mathbb P^1$, and also for every admissible choice of its affine translations and cusp collar, we have
\[
 \omega_{X_c}\simeq\mathcal O_{X_c}(-2S_2),\qquad
 f_{c*}\mathcal O_{X_c}(kS_2)=\mathcal O_B(\lfloor k/4\rfloor p_2)\quad(k\in\mathbb Z).
\]
The latter equality is an equality of subsheaves of meromorphic functions under pullback. The anticanonical section ring is $\mathbb C[U,V]$, with $\deg U=1$, $\deg V=2$. The complete degree-two map is $f_c$ followed by a projective coordinate isomorphism of its base. Every biholomorphism between two such threefolds is over the identity of $B$.
\end{proposition}
\begin{proof}
We compute the line and its evaluation map, including the multiple fibres. A holomorphic function restricts to a constant on every connected compact reduced fibre. To descend functions on a neighbourhood we first use regular fibres; equality then extends from their dense union by the identity theorem on the smooth total space. Away from the multiple fibres, evaluation on local sections descends it to a holomorphic base function. On a finite root cover this follows by evaluation on the zero section; invariance under $s\mapsto\zeta s$ leaves precisely the convergent series in $s^m$. At the cusp, a smooth point of the reduced fibre admits a local section, and constancy on each nearby compact fibre gives the same descent there by the identity theorem. Consequently $f_*\mathcal O_X=\mathcal O_B$.

Put $e_\wedge=d\zeta_1\wedge d\zeta_2$ in relative fibre coordinates. At the finite centre $j$ let $a_1=2,a_2=1$ and $\zeta_j=e^{-2\pi i/m_j}$. The canonical frame $ds_j\wedge e_\wedge$ has character $\chi_j(s)=\zeta_j\det R_j(s)$, with $\chi_j(0)=\zeta_j^{1+a_j}$. Take its holomorphic logarithm $\chi_j=\zeta_j^{1+a_j}\exp\psi_j$, where $\psi_j(0)=0$. The finite group identity says $\sum_{r=0}^{m_j-1}\psi_j(\zeta_j^r s)=0$. Writing $\psi_j=\sum c_ns^n$, we obtain $c_n=0$ for $m_j\mid n$. Therefore
\[
 \varphi_j(s)=\sum_{m_j\nmid n}\frac{c_n}{1-\zeta_j^n}s^n
\]
is convergent and satisfies $\varphi_j(s)-\varphi_j(\zeta_js)=\psi_j(s)$. The invariant coefficient module of the relative canonical frame
$e^{\varphi_j}(ds_j\wedge e_\wedge)\otimes(dt_j)^{-1}$ is exactly
$s_j^{m_j-1-a_j}\mathbb C\{s_j^{m_j}\}$. Notice that $(dt_j)^{-1}$ denotes a pulled-back line-bundle frame; no division by a vanishing differential is made. Since the total-space cover is etale, the evaluation divisor downstairs is $(m_j-1-a_j)S_j$, namely $0$ and $2S_2$.

At the cusp the ordinary form
\[
 \Omega=\frac{dx_1}{x_1}\wedge\frac{dx_2}{x_2}\wedge dt_c
\]
is nowhere zero. Indeed in a maximal unimodular cone chart, with coordinates $z_0,z_1,z_2$, the height-one condition gives $t_c=z_0z_1z_2$ and unimodularity gives $\Omega=\pm dz_0\wedge dz_1\wedge dz_2$. The same equality after inverting coordinates treats the other charts. Integral shears and holomorphic torus multipliers depending on $t_c$ preserve this form: all their extra differential terms contain $dt_c$ and vanish in the wedge. Thus the form descends to the cusp quotient and evaluation has no zero there. This calculation also shows that changing any collar by a fibre translation does not alter the canonical transition factors.

We next determine $\mathfrak K=f_*\omega_{X/B}$. On the simply connected period cover let $h^2=E_6(\tau)$, and put $\Theta=E_4(\tau)^2h/\Delta(\tau)$. The relation $j(\tau)=1728t$ gives local degrees $\tau-\rho\sim s_1$, $\tau-i\sim s_2^2$ and $q=e^{2\pi i\tau}\sim t_c/1728$. The identities $E_4^3=1728t\Delta$ and $E_6^2=1728(t-1)\Delta$, together with the nonvanishing product $\Delta=q\prod_{n\ge1}(1-q^n)^{24}$ on the upper half-plane, give orders $2,1,-1$ for $\Theta$. The same identities exclude every other zero or pole. The square root exists because the zeros of $E_6(\tau)$ have even order and the cover is simply connected. Its signs are determined locally: at $\rho$, $\rho^3=-1$ and $h\ne0$, whereas at $i$, $h$ has order one in $s_2$. The weight transformation laws consequently give
\[
 \Theta(g_1z)=-\Theta(z)/\tau(z),\qquad
 \Theta(g_2z)=\Theta(z)/\tau(z),\qquad \Theta(g_0z)=\Theta(z).
\]
The invariant relative form $\Theta^{-1}e_\wedge$ is a unit in the finite local generator just calculated and has a simple zero in the cusp generator $e_\wedge=(2\pi i)^{-2}\Omega\otimes(dt_c)^{-1}$. It defines an isomorphism $\mathfrak K\simeq\mathcal O_B(p_0)$. The evaluation divisor now yields
\[
 \omega_X\simeq f^*\mathcal O_B(p_0)\otimes f^*\omega_B\otimes\mathcal O_X(2S_2).
\]
In the fixed coordinate $t$, the function $(t-1)^{-1}$ has divisor $p_0-p_2$, while $dt/(t-1)^2$ has divisor $-2p_2$. These explicit maps identify the last display with $\mathcal O_X(-4S_2+2S_2)$.

A meromorphic representative of a section of $\mathcal O_X(kS_2)$ is holomorphic off $S_2$, so the constancy argument makes it a pullback there. On the root cover its value along the zero section is meromorphic, with order at least $-k$, and is invariant under $s\mapsto-is$. Its Laurent series has only powers $s^{4r}$. The necessary and sufficient inequality is $4\operatorname{ord}_{p_2}b+k\ge0$, which proves the direct-image formula, including negative $k$.

Choose the canonical divisor section $u$ of $\mathcal O_X(2S_2)$ and identify this line with $\omega_X^{-1}$ using the preceding maps. Then $u^2=f^*\sigma_{p_2}$. Choose $\ell\in H^0(B,\mathcal O_B(p_2))$ with $\ell(p_2)\ne0$, and let $v=f^*\ell$ in degree two. In degree $m=2r$ the monomials $u^{2r-2j}v^j$, $0\le j\le r$, are the pullbacks of the usual basis $\sigma_{p_2}^{r-j}\ell^j$. In degree $2r+1$, multiplication by $u$ gives the same basis because the allowed base pole order remains $r$. These are all sections by the direct-image formula. Thus $U\mapsto u,V\mapsto v$ is a graded algebra isomorphism, not only a Hilbert-series equality. The two degree-two sections have no common zero and give $f$ in the projective coordinate $[\sigma_{p_2}:\ell]$.

A biholomorphism pulls back the canonical line, its sections and their evaluations; it therefore induces $\alpha\in\operatorname{PGL}_2(\mathbb C)$ with $f'F=\alpha f$. The cusp fibre is reduced with nonnormal points. The two finite fibres have smooth reductions and nilpotence orders exactly three and four in their local equations $s^3=0,s^4=0$. These are intrinsic local-ring properties of the full fibres and distinguish all three. Every other fibre is a smooth reduced torus. Consequently $\alpha$ fixes $\infty,0,1$ individually, and the equation of a fractional linear transformation gives $\alpha=1$.
\end{proof}

\begin{proposition}[A.1.2; Proved]\label{st:A.1.2}
The normalized base change $N_j\times_{D_j}\Delta_j\to\Delta_j$, $t_j=s_j^{m_j}$, is the original smooth torus family $J'_j$. The maps $E_1\times F_1\to S_1$ and $E_2\times F_2\to S_2$ described below are finite etale covers of degrees $9$ and $8$.
\end{proposition}
\begin{proof}
On $J'_j$, the deck generator is
$(s,\zeta)\mapsto(\zeta_js,R_j(s)\zeta+\Pi_c(g_js)v_j/m_j)$. It has exact order $m_j$. Its $r$th power changes the invariant circle coordinate $\gamma$ by $r\gamma(v_j)/m_j$; this is nonintegral for $0<r<m_j$, so the action is free. The quotient map is therefore a local biholomorphism. Its local base function is $s^{m_j}$ and its fibre multiplicity is exactly $m_j$. The map
\[
 J'_j\longrightarrow N_j\times_{D_j}\Delta_j,\qquad y\longmapsto(q_j(y),s_j(y))
\]
is finite, is an isomorphism over $s\ne0$, and has normal source. Locally at the central fibre the base-change ring has equation $z^{m_j}-s^{m_j}=0$. Its normalization is the disjoint union of branches $z=\zeta_j^rs$. The displayed map gives those branches in the covering charts. It is thus the normalization everywhere, with its specified root coordinate.

Put $\rho=e^{\pi i/3}$, $\mu_1=(2-\rho)/3$, $\mu_2=(1-i)/2$, $\kappa_1=\beta_c(z_1)+2\mu_1$, and $\kappa_2=\beta_c(z_2)+3\mu_2$. The fixed-point transformation laws for $\mu$ give the two displayed values directly. We have $\operatorname{Im}\kappa_j=D_c(z_j)<0$. Define
\[
 E_j=\mathbb C/(\mathbb Z+\kappa_j\mathbb Z),\quad
 F_1=\mathbb C/(\mathbb Z(\rho-2)-3\mathbb Z),\quad
 F_2=\mathbb C/(\mathbb Z(i-1)-2\mathbb Z).
\]
Let $e_1=(0,1,-2,0)^t,f_1=(0,0,-3,1)^t,e_2=(0,1,-1,0)^t,f_2=(0,0,-2,1)^t$. Direct solution of the fixed equations and inspection of the columns give
\[
 \Lambda^{A_1}=\mathbb Z\varepsilon+\mathbb Zd,\quad
 \Lambda^{A_2}=\mathbb Z\varepsilon'+\mathbb Zd,\quad
 (A_j-I)\Lambda=\mathbb Ze_j+\mathbb Zf_j.
\]
The image lattices are saturated: they are the integer solutions of $\gamma=0,2u+w+3\delta=0$ and $\gamma=0,u+w+2\delta=0$, respectively. The determinants of the fixed-plus-moving basis matrices are $-3,-2$.

The complex-linear maps $L_1(\xi,\eta)=(\eta,\xi-\eta/3)$ and $L_2(\xi,\eta)=(\eta,\xi-\eta/2)$ carry the product periods to those fixed and moving sublattices. They induce isogenies to the original central tori, with kernel generated respectively by $(2/3,\rho)$ and $(1/2,i)$, of exact orders three and two. For example $\widehat u=2d/3+e_1-2f_1/3=d/2+e_2-f_2/2$ gives both kernel representatives.

The affine deck maps lift to
\[
 (\xi,\eta)\longmapsto(\xi+\kappa_1/3,(\rho-1)\eta),\qquad
 (\xi,\eta)\longmapsto(\xi-\kappa_2/4,-i\eta).
\]
Substitution in $L_j$ gives exactly $R_jL_j+\Pi_cv_j/m_j$. They commute with their kernel translations: the moving differences are periods, since $(\rho-2)\rho=-1-\rho=-(\rho-2)-3$ can equally be checked in the lattice basis, and $(-i-1)i=1-i=-(i-1)$. More directly these assertions follow by the two-dimensional period matrices of the indicated rotations. A nonidentity element of either resulting product group translates the $E_j$ coordinate by $2a/3+b\kappa_1/3$ or $a/2-b\kappa_2/4$, respectively. Real independence of $1,\kappa_j$ proves that this translation vanishes only for the identity element. Hence the groups $C_3\times C_3$ and $C_2\times C_4$ act freely and have quotient precisely $S_j$.
\end{proof}

\subsection{A.2. The integral obstruction}

\begin{proposition}[A.2.1; Proved]\label{st:A.2.1}
The rational centralizer of $(A_1,A_2)$ is $\mathbb QI+\mathbb QE$. Every biholomorphism $X_c\to X_{c'}$ has $c'-c=n\in2\mathbb Z$, complex-linear fibre part $I_2$, and source-to-target lattice map $I-nE$. In the full fixed lattices,
\[
 \Lambda^{A_1}/N_1\Lambda\simeq\mathbb Z/3,\qquad
 \Lambda^{A_2}/N_2\Lambda\simeq\mathbb Z/4\oplus\mathbb Z/2,
\]
with $d$ generating the second factor of the latter group and vanishing in the former.
\end{proposition}
\begin{proof}
The common invariant vector space is $\mathbb Qd$, and the common invariant covector space is $\mathbb Q\gamma$; this follows by intersecting the fixed bases in A.1.2 and the rows $\langle\gamma,2u+w+3\delta\rangle$, $\langle\gamma,u+w+2\delta\rangle$. A commuting endomorphism preserves $\mathbb Qd\subset\ker\gamma$. Its induced matrix on $\ker\gamma/\mathbb Qd$ commutes with
$Q_1=\left(\begin{smallmatrix}0&1\\-1&-1\end{smallmatrix}\right)$ and $Q_2=\left(\begin{smallmatrix}0&-1\\1&0\end{smallmatrix}\right)$. Commutation with $Q_2$ makes it $\left(\begin{smallmatrix}b&r\\-r&b\end{smallmatrix}\right)$, and the $Q_1$ equation sets $r=0$. Write the remaining columns as
$H\widehat\gamma=(a,x,y,k)^t$, $H\widehat u=b\widehat u+u_0d$, $H\widehat w=b\widehat w+w_0d$, $Hd=d_0d$. The equations on $\widehat u,\widehat w$ give $d_0-w_0=b+u_0$, $u_0=2w_0$ and $u_0=w_0$, hence $u_0=w_0=0,d_0=b$. The first-column equations give $x=2(a-b),y=-4(a-b),x=-y$. Thus $a=b,x=y=0$, proving the assertion. Conversely $A_jE=EA_j=E$.

A fibre biholomorphism lifts to an affine map on $\mathbb C^2$: the derivative of a lifted map is lattice periodic, descends to a holomorphic function on a compact torus, and is constant. Its lattice map is locally constant and commutes with both monodromies. Write it as $H=aI+bE$; integrality and $\det H=a^4=1$ give $a=\pm1,b\in\mathbb Z$. The equation $L\Pi_c=\Pi_{c'}H$, in columns three and four, gives $L=aI_2$; column one gives $a(c'-c)+b=0$. Thus $n=c'-c\in\mathbb Z$.

The normalized base-change construction in A.1.2 is functorial for biholomorphisms over the base: the base-changed isomorphism transports integral elements of the total quotient ring to integral elements, and its inverse does the same. Its lift is over the identity of the root disc. It commutes with the prescribed deck generator because both compositions lift the same map $(x,s)\mapsto(F(x),\zeta_js)$ on the unnormalized base change. This proves exact commutation, rather than commutation after an unchosen power of the generator.

Lift the translation of the resulting affine torus map holomorphically to a vector $B_j(s)$ on the simply connected root disc. Its deck equation, expressed in the source periods, is
\[
 B_j(g_jz)-R_j(z)B_j(z)=\Pi_c(g_jz)
 \left(\frac{(1-a)v_j+n\gamma(v_j)d}{m_j}+\ell_j\right),\qquad\ell_j\in\Lambda.
\]
The integral vector is constant by continuity and real invertibility of the period matrix. Iterating the equation $m_j$ times gives
\[
 N_j\ell_j=-(1-a)v_j-n\gamma(v_j)d.
\]
Exact multiplication yields
\[
 N_1=\begin{pmatrix}3&0&0&0\\6&0&0&0\\-12&0&0&0\\0&2&1&3\end{pmatrix},\quad
 N_2=\begin{pmatrix}4&0&0&0\\12&0&0&0\\-12&0&0&0\\0&2&2&4\end{pmatrix}.
\]
In fixed bases their images are $3\mathbb Z\varepsilon+\mathbb Zd$ and $4\mathbb Z\varepsilon'+2\mathbb Zd$. These explicit generators prove the asserted quotient groups. Applying $\gamma$ to the norm equation excludes $a=-1$ already at $m_1=3$, since $3\gamma(\ell_1)=-2$. For $a=1$, the second norm equation says exactly $2\mid n$; the first imposes no condition. Program \texttt{code/ab\_exact.py} verifies the matrices, Smith forms, centralizer and membership independently by Hermite-form comparisons. All its computations are over $\mathbb Z$ or $\mathbb Q$.
\end{proof}

\subsection{A.3. Extension and coherent descent}

\begin{proposition}[A.3.1; Proved]\label{st:A.3.1}
For each $n\in2\mathbb Z$ there is a biholomorphism $F_{n,c}:X_c\to X_{c+n}$, holomorphic jointly in $c$, with $F_{2h,c+2k}F_{2k,c}=F_{2(h+k),c}$. The quotient family over $0<|w|<e^{-\pi M}$, $w=e^{-\pi ic}$, is proper and holomorphic, with pairwise non-biholomorphic fibres.
\end{proposition}
\begin{proof}
Integer shifts satisfy $\Pi_{c+n}=\Pi_c(I+nE)$, so the middle lattices agree. Set
\[
 r_1=-\tfrac n3(\widehat u+\widehat w),\quad
 r_2=\tfrac n4(\widehat u+\widehat w),\quad
 \ell_1=-n\widehat w,\quad \ell_2=\tfrac n2\widehat u,\quad b_j=\Pi_cr_j.
\]
The four-entry equations are $(I-A_1)r_1=(0,0,-n,n/3)^t$ and $(I-A_2)r_2=(0,n/2,0,-n/4)^t$. Consequently $B_j=b_j$ solves the preceding deck equation for $a=1$. The map $(z,\zeta)\mapsto(z,\zeta+b_j(z))$ extends on the entire root family and descends, together with its inverse, to each finite filling.

Its middle translation on the punctured overlap is
\[
 a_j=b_j+n\gamma(v_j)\frac{\log s_j}{2\pi i}e.
\]
This follows from the actual overlap $h_{j,c}(\zeta)=\zeta+\sigma_{j,c}$, where $\sigma_{j,c}=(\log s_j)/(2\pi i)\,\Pi_cv_j$: precisely $\sigma_{j,c+n}-\sigma_{j,c}=n\gamma(v_j)(\log s_j)/(2\pi i)e$. Changes of logarithm add periods. Since the chosen generator changes $(\log s_j)/(2\pi i)$ by $-1/m_j$, we have $a_j(g_jz)=R_ja_j(z)+\Pi_c(g_jz)\ell_j$.

On the universal cover of $B^\circ$, extend $\ell_1,\ell_2$ by the integral cocycle law $\ell_{gh}=\ell_g+A_g\ell_h$. Freeness of $\pi_1(B^\circ)$ makes this well defined. At $g_0=(g_1g_2)^{-1}$ it gives
\[
 \ell_0=-A_0(\ell_1+A_1\ell_2)=\tfrac{3n}2\widehat w-\tfrac n2d.
\]
Put $q_n=(-3n/2,n/2)^t$ and $a_0=sq_n$. Then $a_0(g_0z)-a_0(z)=\Pi_c\ell_0$. In cusp torus coordinates its map is $(x,t_c)\mapsto(t_c^{q_n}x,t_c)$. The integral linear map $(y,h)\mapsto(y+hq_n,h)$ of the fan lattice translates each height-one triangle by an integer vector. It therefore permutes every cone of the prescribed periodic triangulation. The map and its inverse extend across all strata, including both triple points and all three double curves after quotient. Moreover the cusp deck multipliers are identical at $c$ and $c+n$, since $C_{c+n}\lambda-C_c\lambda=(0,n\lambda_1)^t$ has integral exponential. Thus the extended map descends with all its nonconstant units retained.

We supply the global analytic correction. Quotient the trivial affine vector bundle on the universal cover by
$(z,y)\mapsto(gz,R_g(z)y+\Pi_c(gz)\ell_g)$. Its action law is exactly the cocycle law by period covariance. Specify at the finite points sections of the form $a_j+\eta_j$, with $\eta_j$ holomorphic equivariant on the full root disc; at the cusp specify $a_0+\eta_0(t_c)$ with $\eta_0$ holomorphic on the full disc. These prescriptions define a locally nonempty torsor under a rank-two locally free sheaf $\mathcal V$ on $B$.

Projection to the first coordinate gives
\[
 0\longrightarrow\mathcal O_Be\longrightarrow\mathcal V
 \longrightarrow\mathcal O_B(-p_0)\longrightarrow0.
\]
Indeed its factors are $-1/\tau,1/\tau,1$. Multiplication by $\Theta$ identifies this first-coordinate sheaf with holomorphic functions at ordinary and finite points and functions vanishing at $p_0$: the finite vanishing orders are $2,1$, and the cusp pole is one. At a finite stalk, any equivariant first component has an equivariant vector lift by averaging a preliminary lift $(u,0)$ over the finite group. This proves stalkwise surjectivity and local freeness, including the off-diagonal entries of $R_j$.

Both $H^1(\mathbb P^1,\mathcal O)$ and $H^1(\mathbb P^1,\mathcal O(-1))$ vanish. To make the analytic gluing explicit, split a smooth cocycle with a partition of unity, solve $\bar\partial$ locally by the Cauchy integral, and compare the two resulting solutions on an annulus. A Laurent series $\sum a_nt^n$ is the difference of its nonnegative and negative parts for $\mathcal O$; for $\mathcal O(-1)$ use the second frame $t^{-1}$, so the negative terms become $t^{-1}(-\sum_{n<0}a_nt^{n+1})$. Both series converge on their respective discs and leave no exponent unsolved. This kills the cocycle for any refinement. The exact sequence gives $H^1(B,\mathcal V)=0$. Project the torsor cocycle to $\mathcal O(-1)$, solve it, lift the solutions, and solve the remaining scalar cocycle. We obtain one global section $a$, with $a=a_j+\eta_j$ at finite points and $a=sq_n+\eta_0$ at the cusp.

The three actual maps are now
\[
 \zeta\mapsto\zeta+a,\qquad
 \zeta\mapsto\zeta+b_j+\eta_j,\qquad
 (x,t_c)\mapsto(e^{2\pi i\eta_0(t_c)}t_c^{q_n}x,t_c).
\]
The finite overlap identity is the equality of sums
$\sigma_{j,c}+a=\sigma_{j,c+n}+b_j+\eta_j$. The cusp overlap identity is $e^{2\pi ia}=t_c^{q_n}e^{2\pi i\eta_0}$. The holomorphic-unit factors act on every toric chart. The local inverses satisfy the same identities; hence the maps glue to a biholomorphism.

For coherence, perform the additive solution only for $n=2$ and multiply that solution by $k$ for $n=2k$. The cocycles have $\gamma(\ell_g)=0$, so their affine terms $\Pi_c\ell_g$ are independent of $c$. The reference translations, $\mathcal V$, and the correction can all be chosen independent of $c$. Adding the translations and multiplying the toric factors proves the stated cocycle identity exactly.

Finally the original gluing constructs a holomorphic family over any parameter disc compactly contained in $\mathcal P$. In ordinary and finite charts, real invertibility of $\Pi_c$ is uniform on compact parameter sets, so a fixed compact real parallelepiped proves properness and the lattice quotient is holomorphic. At the cusp write $Z=sB_0+C(t_c,c)$ with $B_0=\left(\begin{smallmatrix}0&1\\-1&0\end{smallmatrix}\right)$. On a compact parameter disc $C$ is uniformly bounded; shrink the cusp radius until
$\|(-2\pi\operatorname{Im}C)/\log|t_c|\|<1/2$. The effective logarithmic lattice is then within $1/2$ of the unimodular $B_0$, uniformly. Reduce its logarithmic valuation modulo the effective logarithmic lattice to a bounded parallelogram. Choose a height-one triangle containing that reduced valuation. Its barycentric coordinates are nonnegative and sum to one; the corresponding three unimodular cone characters therefore have modulus at most one. This puts the point in the closed unit polydisc of that cone. Only finitely many triangles meet the fixed bounded parallelogram, uniformly on the parameter compactum. Boundary orbits have the same representatives by translating their height-one components to these triangles. Thus every orbit has a representative in this finite union of cone polydiscs. Those closed polydiscs times a compact parameter set form a compact surjective representative set. This proves properness at the cusp and proper discontinuity by the same uniform lower bound on nonzero lattice translations. The quotient charts are open subsets of a fixed smooth toric variety times the parameter disc, so projection to the parameter is a submersion. The explicit overlap maps are jointly holomorphic. These local families agree after shrinking cusp radii, and glue over $\mathcal P$. The coherent $2\mathbb Z$ action is properly discontinuous because its action on $\mathcal P$ is. Local inverse logarithms for $w=e^{-\pi ic}$ identify the quotient family with these proper local families. A.2.1 gives pairwise non-biholomorphic fibres.
\end{proof}

\subsection{A.4. Albanese invariants and the derivative of the deformation}

\begin{proposition}[A.4.1; Proved]\label{st:A.4.1}
With $\kappa_j(c)=\kappa_j^{\rm part}+c$, the Albanese elliptic curves of the reduced finite fibres are
\[
 \operatorname{Alb}(S_1)=\mathbb C/(\mathbb Z+\kappa_1(c)\mathbb Z),\qquad
 \operatorname{Alb}(S_2)=\mathbb C/(\mathbb Z+\tfrac12\kappa_2(c)\mathbb Z).
\]
Writing $J(q)=q^{-1}+744+196884q+\cdots$, their invariants are
\[
 j_1(w)=J(e^{-2\pi i\kappa_1^{\rm part}}w^2),\qquad
 j_2(w)=J(e^{-\pi i\kappa_2^{\rm part}}w).
\]
These are convergent formulas because the displayed nome moduli are below one. Outside the explicitly defined countable exceptional set in the proof, equality of both Albanese invariants forces $c'-c\in2\mathbb Z$. Thus the full generic Albanese cross-check of Lead 3 holds independently of the centralizer calculation.
\end{proposition}
\begin{proof}
In the product covers of A.1.2 the invariant holomorphic one-forms are exactly the multiples of $d\xi$: the moving factor is multiplied by a primitive third or fourth root. Integration thus identifies the Albanese with the quotient of the fixed line by the fixed-coordinate projections of every deck translation. These are, respectively,
\[
 \mathbb Z+\kappa_1\mathbb Z+\tfrac23\mathbb Z+\tfrac{\kappa_1}3\mathbb Z
   =\tfrac13\mathbb Z+\tfrac{\kappa_1}3\mathbb Z,
\quad
 \mathbb Z+\kappa_2\mathbb Z+\tfrac12\mathbb Z+\tfrac{\kappa_2}4\mathbb Z
   =\tfrac12\mathbb Z+\tfrac{\kappa_2}4\mathbb Z.
\]
Multiplication by three and two gives the stated maps of complex tori. In particular the Albanese computation includes the kernel of the product isogeny and the affine quotient, rather than just its fixed eigenspace. Using the positive-imaginary representatives $-\kappa_1$ and $-\kappa_2/2$ in the standard modular $j$-series gives the formulas. At sufficiently small nome $J$ is injective: $1/J(q)=q+O(q^2)$ has nonzero derivative at zero. Thus $j_2$ alone distinguishes parameters within a sufficiently small punctured subdisc when both parameters lie there. For the full generic assertion put $x=-\kappa_2(c)/2$ and $d_*=\kappa_1^{\rm part}-\kappa_2^{\rm part}$, so $-\kappa_1(c)=2x-d_*$. The independently certified interval in \texttt{cd\_special\_intervals.py} gives $-0.2<\operatorname{Im}d_*<-0.1$; in particular it is nonzero. Equality of the two elliptic invariants at $c,c'$ gives matrices $M,N\in\operatorname{PSL}_2(\mathbb Z)$ with $x'=M(x)$ and
\[
 N(2x-d_*)=2M(x)-d_*.
\]
Define the exceptional set to be the $c\in\mathcal P$ obtained from roots of these equations for all pairs for which the two rational functions are not identical. Each equation has at most two roots after clearing denominators, so this is a countable, explicitly specified set. For a pair giving an identity, let $C(x)=2x-d_*$. Then $N=CMC^{-1}$. The real matrix $M$ preserves both the real circle $\mathbb R\cup\{\infty\}$ and $C^{-1}(\mathbb R\cup\{\infty\})$, the parallel line with imaginary part $\operatorname{Im}d_*/2$ together with infinity. Their unique intersection is infinity, so $M$ fixes infinity. An element of $\operatorname{PSL}_2(\mathbb Z)$ fixing infinity is $x\mapsto x+b$, $b\in\mathbb Z$, and then $N(y)=y+2b$. It follows that $c'-c=-2b$. This proves the generic necessity with its actual exceptional set; A.2.1 proves the necessity also on that set.
\end{proof}

\begin{proposition}[A.4.2; Proved]\label{st:A.4.2}
For every $c\in\mathcal P$, the Kodaira--Spencer class of $\partial/\partial c$ in $H^1(X_c,T_{X_c})$ is nonzero. On any fixed ordinary fibre, put $T=\operatorname{Im}\tau$, $m=\operatorname{Im}\mu$, $a=m/T$, $D=\operatorname{Im}\beta_c-6m^2/T$. In the convention $\mathrm{KS}(\partial_c)=[\bar\partial h]$ for a smooth $(1,0)$ lift $h$, its restricted torus-deformation representative is
\[
 -\frac{1}{2iD}\,\partial_{\zeta_2}\otimes
       (d\bar\zeta_2-a\,d\bar\zeta_1).
\]
\end{proposition}
\begin{proof}
Retain the same real marking $r\in\mathbb R^4$ and set $\zeta=\Pi_cr$. The two imaginary-part equations give
$r_1=(\operatorname{Im}\zeta_2-a\operatorname{Im}\zeta_1)/D$. Thus at $c+\delta$ the marked real-linear identification is
\[
 w=\zeta+\frac{\delta}{2iD}P(\zeta-\bar\zeta),\qquad
 P=\begin{pmatrix}0&0\\-a&1\end{pmatrix},\qquad P^2=P.
\]
It is exactly $\Pi_{c+\delta}\Pi_c^{-1}$ as a real map, so it sends every marked lattice vector to the corresponding vector. Its determinant is $(D+\operatorname{Im}\delta)/D>0$. The transported $(0,1)$ graph coefficient is $\nu_\delta=\delta P/(2iD+\delta)$. The derivative at zero is $P/(2iD)$. Equivalently the smooth lift of $\partial_c$ in moving coordinates has vertical part $r_1\partial_{w_2}$ and $\bar\partial r_1=-(d\bar\zeta_2-a d\bar\zeta_1)/(2iD)$, giving the stated sign.

A nonzero constant tangent-valued $(0,1)$ form on a compact torus cannot be $\bar\partial$ of a global vector field. Write that field in the invariant holomorphic frame and integrate each derivative of its periodic coefficients against invariant probability volume. Every such derivative has integral zero, whereas the coefficient of $\partial_{\zeta_2}\otimes d\bar\zeta_2$ above has nonzero integral. This proves nonvanishing for the fibre deformation.

We explain its transfer to the threefold. Over the dual numbers $A=\mathbb C[\epsilon]/(\epsilon^2)$ at $c$, the chart proof of A.1.1 gives the relative line isomorphism $\omega^{-2}_{X_A/A}=f_A^*\mathcal O_B(p_2)$. Extra parameter differentials disappear in relative forms and all determinant characters are independent of $c$. The two pulled-back sections form an $A$-basis of $H^0(X_A,\omega^{-2})$: reduce a section modulo $\epsilon$, subtract its linear combination of the two lifts, and use the exact sequence $0\to L\xrightarrow{\epsilon}L_A\to L\to0$ to write the remainder uniquely as $\epsilon$ times another such combination. The same argument proves independence.

A first-order trivialization would preserve this complete system and give a base automorphism reducing to the identity. It fixes all three special values even over $A$. On the finite etale charts the relative critical ideal is $(s_j^{m_j-1})$, and $t_j=s_j^{m_j}$ is zero on that scheme; its constant parameter ring $A$ injects. At a cusp triple chart, the critical ideal of $t_c=z_0z_1z_2$ is $(z_0z_1,z_0z_2,z_1z_2)$ and has the same property. There are no other critical values. Therefore an infinitesimal displacement of a special value would remain nonzero in its critical chart, a contradiction. Writing a base infinitesimal automorphism as $t+\epsilon(b_0+b_1t+b_2t^2)$, fixing $0,1,\infty$ gives $b_0=b_1=b_2=0$. The trivialization would consequently trivialize the displayed nonzero torus deformation. Finally a zero Kodaira--Spencer class produces such a first-order trivialization by subtracting from local product charts the vector fields whose coboundary is the differentiated transition cocycle. This proves the claim.
\end{proof}

\begin{proposition}[A.4.3; Proved]\label{st:A.4.3}
The non-isotriviality argument of Viaclovsky, \cite[Proposition 4.7]{VL26}, follows for the present family from A.4.2 and A.1.1. His published countability argument proves local non-isotriviality and positive Kuranishi dimension; the computation in A.4.2 additionally proves nonzero derivative at every admissible parameter. The arithmetic method of Yang, \cite[Lemmas 6.6--6.9 and Proposition 6.15]{Y26}, specializes to the same centralizer, finite-norm and collar-gluing operations after replacing his displayed matrices by $A_1,A_2$; those replacements are evaluated in A.2.1 and A.3.1.
\end{proposition}
\begin{proof}
On a fixed ordinary fibre, the three last periods remain fixed and span $\mathbb C^2$, whereas the first changes by $\delta e$. A marked complex-linear identification fixing those three periods must be the identity, so the marked period map is nonconstant. The set of integral markings of one torus is countable. A holomorphic nonconstant map from a disc has uncountable image; hence the unmarked fibre family is locally non-isotrivial. A.1.1 transfers this to total spaces. This is exactly the implication used in the cited Proposition 4.7, whose statement and proof we have read. A nonconstant classifying map gives positive dimension of its local target. It does not by itself determine the derivative at the base point, as the explicit holomorphic map $z\mapsto z^2$ has nonconstant image and derivative zero at zero. A.4.2 computes the derivative here.

Yang's cited lemmas use an integral commuting lattice map, an affine fibre lift, normalization of the root base change, and a norm-lattice defect. We have exhibited the respective maps and norm equations for the present matrices. His Proposition 6.15 obtains $6/r$ from the first-column scaling and its order-six norm image. Here the first-column scaling is one and the two computed norm images impose parity, yielding two. The common proof operations transfer through these explicitly stated matrix substitutions; no conclusion about one family is inferred merely from the label of the other. His nonzero absolute Kodaira--Spencer conclusion is obtained in the present setting by the explicit computation and intrinsic-pencil transfer above.
\end{proof}

\subsection{A.5. The assembled theorem}

\begin{theorem}[A.5.1; Proved]\label{st:A.5.1}
Fix the original cusp-regular period solution, the clockwise peripheral marking, the prescribed periodic triangular cusp fan, the affine vectors $(\varepsilon,-\varepsilon')$, and the untwisted cusp collar. For every $c,c'\in\mathcal P$, the constructed compact complex threefolds satisfy
\[
 X_c\simeq X_{c'}\quad\Longleftrightarrow\quad c'-c\in2\mathbb Z.
\]
Every biholomorphism is over the identity of the original $\mathbb P^1$. The proper holomorphic family of their isomorphism classes has parameter $0<|w|<e^{-\pi M}$ and nonzero Kodaira--Spencer class at every point.
\end{theorem}
\begin{proof}
A.1.1 makes the base intrinsic. A.2.1 gives the necessary parity and excludes the negative linear part. A.3.1 supplies the inverse maps, all collar identities, and descent of the jointly holomorphic family. A.4.2 proves its infinitesimal effectiveness. Finally $|e^{-\pi ic}|=e^{\pi\operatorname{Im}c}$ and equality of exponentials is exactly difference in $2\mathbb Z$, so the coordinate has precisely the stated image and fibres.
\end{proof}

\begin{theorem}[A.5.2; Proved]\label{st:A.5.2}
For every admissible $c$, the tangent cohomology of the actual compact
member is
\[
 \bigl(h^0(X_c,T_{X_c}),h^1(X_c,T_{X_c}),h^2(X_c,T_{X_c}),
 h^3(X_c,T_{X_c})\bigr)=(1,1,1,0).
\]
Its Kuranishi germ is a smooth reduced disc. The family of A.5.1 is
miniversal at every parameter: its Kodaira--Spencer map is an isomorphism,
and every sufficiently small deformation is induced from this family.
\end{theorem}
\begin{proof}
We give the calculation, including the singular-fibre maps. The
corresponding calculation appears in the historical first version
\cite[Section 6, especially Proposition 6.10, Lemma 6.11,
Theorem 6.12 and Proposition 6.16]{DGKX26v1}; those deformation
statements are omitted from version two. We use the verified maps and
cohomology calculation below, rather than infer their validity from the
earlier theorem label.

\emph{The base and its line bundles.}
Write $X=X_c$ and form the analytic root curve
\[
 \mathfrak B=\sqrt[3]{(B,p_1)}\mathbin{\times_B}\sqrt[4]{(B,p_2)},
 \qquad \rho:\mathfrak B\longrightarrow B.
\]
A root object is a line bundle with section whose indicated power is the
pullback of the section cutting out $p_j$. The actual equalities
$f^*p_1=3S_1$ and $f^*p_2=4S_2$ supply these objects as
$(\mathcal O_X(S_j),1_{S_j})$, and hence supply a morphism
$\bar f:X\to\mathfrak B$ with $\rho\bar f=f$.
The explicit weighted quotient
\[
 \mathfrak B\simeq[(\mathbb C^2\setminus\{0\})/\mathbb C^*],
 \qquad \lambda(x,y)=(\lambda^3x,\lambda^4y),
 \qquad t=\frac{y^3}{y^3-x^4},
\]
identifies the roots with $(\mathcal O(4),y)$ and
$(\mathcal O(3),x)$ and the ordinary point $p_0$ with
$y^3-x^4=0$. In particular $\rho^*\mathcal O_B(1)=\mathcal O(12)$.
This isomorphism follows on each root chart by the displayed root
sections; away from their zeros both sides have the same ordinary coarse
coordinate, so the chart isomorphisms glue.

For completeness the analytic line-bundle computation can be made on
the coarse curve. Removing the two stabilizer characters from any line
bundle makes it descend to $B$: on each root chart, averaging a local
frame whose central value is nonzero gives an invariant frame still
nonzero at the centre, and its invariant coefficients are holomorphic
in $s_j^{m_j}$. These frames give the descended line and its transition
functions. Consequently its Picard group has
generators $H,R_1,R_2$ and relations $3R_1=H=4R_2$.
Thus it is freely generated by $R_1-R_2$, denoted $\mathcal O(1)$;
$H=12\mathcal O(1)$, $R_1=4\mathcal O(1)$ and
$R_2=3\mathcal O(1)$. The chart descent just used is exact because
invariants of a finite group over $\mathbb C$ are an exact functor.
With the clockwise root generators, a local invariant section of
$\mathcal O(k)$ has vanishing exponents
\[
 r_3(k)\in\{0,1,2\},\quad r_3(k)\equiv k\pmod3,
 \qquad
 r_4(k)\in\{0,1,2,3\},\quad r_4(k)\equiv-k\pmod4.
\]
Dividing by those local powers gives the precise coarse direct image
\begin{equation}\label{eq:ab-root-degree}
 \rho_*\mathcal O(k)=\mathcal O_B\bigl(d(k)\bigr),\qquad
 d(k)=\frac{k}{12}-\frac{r_3(k)}3-\frac{r_4(k)}4.
\end{equation}
Indeed $k-4r_3-3r_4=12d(k)$ identifies the line bundle with
$\rho^*\mathcal O_B(d(k))\otimes R_1^{r_3}\otimes R_2^{r_4}$.
The canonical root sections generate the invariant direct image of
the last two factors; division by their local powers is the displayed
line-bundle isomorphism on the coarse charts. Finite invariants also give
$R^q\rho_*=0$ for $q>0$, as seen on each Stein quotient chart.
For $k=0$ the degree is zero; for
$k=1,-1,-2,-5,-6$ it is $-1$. It follows directly from the
cohomology of $\mathbb P^1$ that these five line bundles are acyclic,
whereas $H^0(\mathcal O)=\mathbb C$ and $H^1(\mathcal O)=0$.
The local equation $dt_j=m_js_j^{m_j-1}ds_j$ gives
\[
 K_{\mathfrak B}=\mathcal O(-24+2\cdot4+3\cdot3)
 =\mathcal O(-7),\qquad
 T_{\mathfrak B}(-p_0)=\mathcal O(-5).
\]

\emph{The actual base changes and cusp cohomology.}
On a finite root disc $\Delta_{s_j}$ the local presentations are
$[J'_j/\mu_{m_j}]\to[\Delta_{s_j}/\mu_{m_j}]$.
The fibre-product identity
\[
 [J'_j/\mu_{m_j}]
 \mathbin{\times_{[\Delta_{s_j}/\mu_{m_j}]}}\Delta_{s_j}
 \simeq J'_j
\]
is induced by the quotient torsor and its equivariant map to the disc;
its inverse sends a torsor with a chosen root-chart trivialization to
the corresponding point of $J'_j$. The action on $J'_j$ is free by
A.1.2. Thus this is an etale atlas base change and is exactly the
smooth torus family, with its given deck generator. Over the cusp the
base change is the original cusp family. Locally it is
$q=z_0z_1z_2$; the homomorphism
$\mathbb C\{q\}\to\mathbb C\{z_0,z_1,z_2\}$ is flat because its
target is torsion-free over the discrete valuation ring.
These charts prove that $\bar f$ is representable, proper and flat,
with connected Gorenstein fibres of dimension two. Properness on the
cusp is the compact representative construction of A.3.1, and on the
finite charts it is the compact torus quotient. Also
$\bar f^*p_0=W$ as reduced Cartier divisors.

Let $\nu:E\to W$ be the normalization, $E=\mathrm{dP}_6$, let
$C_h,C_v,C_d$ be the three double-curve normalizations, and let
$P_+,P_-$ be the two triple points. Each $C_e$ is $\mathbb P^1$.
Write $\rho_e^0,\rho_e^1:C_e\to E$ for its two paired boundary maps
and $j_e=\nu\rho_e^0=\nu\rho_e^1$. The sheaf sequence is
\begin{equation}\label{eq:ab-cusp-resolution}
 0\longrightarrow\mathcal O_W\longrightarrow\nu_*\mathcal O_E
 \xrightarrow{\delta^0}\bigoplus_{e=h,v,d}j_{e*}\mathcal O_{C_e}
 \xrightarrow{\delta^1}\mathbb C_{P_+}\oplus\mathbb C_{P_-}
 \longrightarrow0,
\end{equation}
where $\delta^0(s)_e=(\rho_e^1)^*s-(\rho_e^0)^*s$ and
$\delta^1$ evaluates $t_h-t_v+t_d$ at the two ordered triple points.
Its exactness is local: for $xyz=0$ it is the alternating restriction
complex of the three coordinate planes, their pairwise lines and their
common point. Decomposing convergent power series by the variables
occurring in each monomial gives the augmented simplex complex, which
is exact. The one- and two-branch cases are the same calculation with
fewer planes. This verifies exactness for the actual boundary maps,
including the torus multipliers in their gluing.

All positive-degree structure-sheaf cohomology of $E,C_e$ vanishes.
Consequently \eqref{eq:ab-cusp-resolution} computes cohomology by
\[
 \mathbb C\xrightarrow{0}\mathbb C^3
 \xrightarrow{\left(\begin{smallmatrix}1&-1&1\\1&-1&1\end{smallmatrix}\right)}
 \mathbb C^2.
\]
It gives $(h^0,h^1,h^2)(\mathcal O_W)=(1,2,1)$.
We also need its multiplication. The branch-intersection resolution
is the augmented semisimplicial sheaf algebra with its restriction face
maps; its local exactness makes the augmentation a multiplicative
quasi-isomorphism. Its Alexander--Whitney product therefore computes
the sheaf-cohomology cup product. In the ordered triangles
$[0,e_1,e_2]$ and $[e_1,e_2,e_1+e_2]$ it is
\[
 \alpha\smile\beta=(\alpha_h\beta_d,\alpha_d\beta_h),
 \qquad [\alpha]\smile[\beta]
 =\alpha_h\beta_d-\alpha_d\beta_h
 \quad\hbox{in }\mathbb C^2/\mathbb C(1,1).
\]
The cocycles $(1,1,0)$ and $(0,1,1)$ pair to one. Thus
$\bigwedge^2H^1(\mathcal O_W)\to H^2(\mathcal O_W)$ is an
isomorphism, including its product rather than just its dimensions.

\emph{Vertical fields and the direct images.}
Set $\mathcal T=T_X(-\log W)$ and
$\mathcal V=\ker(d\bar f^{\log})$.
The vertical transitions on the core are precisely $R_g(z)$.
Finite collar translations have vertical differential the identity,
and at the cusp the frame $\partial_{\zeta_i}$ corresponds to
$2\pi i x_i\partial_{x_i}$. The cusp transformations
$x\mapsto u(q)q^{B_0\lambda}x$ fix the coefficients of these
relative logarithmic fields. These exact identifications construct
a bundle $\mathcal E$ on $\mathfrak B$ with
$\mathcal V=\bar f^*\mathcal E$.
The fixed vector $e=(0,1)^t$ is nonzero in every logarithmic frame,
so there is an exact sequence
\[
 0\longrightarrow\mathcal O\xrightarrow{e}\mathcal E
 \longrightarrow L\longrightarrow0.
\]
The factors of $L$ are $-1/\tau$ and $1/\tau$. Therefore the usual
weight-twelve modular discriminant makes $1/\Delta(\tau)$ a
meromorphic section of $L^{12}$. It has no finite zero or pole,
and its cusp pole is simple: $\Delta(\tau)=q\,u(q)$ with
$u(0)\ne0$. Hence $L^{12}=\mathcal O(-p_0)=\mathcal O(-12)$.
The Picard calculation gives $L=\mathcal O(-1)$, and the extension
splits since $H^1(\mathcal O(1))=0$. Thus
\[
 \mathcal E\simeq\mathcal O\oplus\mathcal O(-1).
\]
There is also the exact sequence
\begin{equation}\label{eq:ab-log-tangent}
 0\longrightarrow\bar f^*\mathcal E\longrightarrow\mathcal T
 \longrightarrow\bar f^*\mathcal O(-5)\longrightarrow0.
\end{equation}
On the smooth root charts it is the relative tangent sequence. At the
cusp its last map sends $\sum a_i z_i\partial_{z_i}$ to
$(\sum a_i)q\partial_q$, so it is onto and its kernel is the
two height-zero lattice directions already used above. Taking its
determinant and using $W=\bar f^*p_0$ gives the actual relative
dualizing line
\begin{equation}\label{eq:ab-relative-dualizing}
 \omega_{X/\mathfrak B}=\bar f^*(\det\mathcal E)^{-1}.
\end{equation}
In particular this computation gives
$K_X=\bar f^*\mathcal O(-6)=\mathcal O_X(-2S_2)$, with the last
identification supplied by $\bar f^*R_2=\mathcal O_X(S_2)$;
it agrees through this actual pullback with A.1.1.

The fibres of $\bar f$ all have structure-sheaf cohomology dimensions
$(1,2,1)$: on ordinary and finite root charts they are tori, and at
the cusp we have just computed them. Analytic proper flat base change
on these complex-space charts gives locally free $R^q\bar f_*\mathcal O_X$
with these fibres. Its canonical comparison maps agree on the atlas
double overlaps, so they descend to $\mathfrak B$. The unit gives
$R^0\bar f_*\mathcal O_X=\mathcal O$.
The source charts are smooth and the fibres are Cartier divisors,
so relative analytic duality applies in its Gorenstein form. The
relative trace pairing, compatible on the same overlaps, and
\eqref{eq:ab-relative-dualizing} give
\[
 (R^2\bar f_*\mathcal O_X)^\vee
 =\bar f_*\omega_{X/\mathfrak B}
 =(\det\mathcal E)^{-1},
 \qquad R^2\bar f_*\mathcal O_X=\mathcal O(-1).
\]
Here local freeness removes the higher sheaf-Ext terms in the duality
spectral sequence. This is precisely the pointwise base-change and
top-degree duality application spelled out in the cited Lemmas 6.8--6.9.

Put $G=R^1\bar f_*\mathcal O_X$. Cup product and the verified fibre
comparisons show $\det G=R^2\bar f_*\mathcal O_X=\mathcal O(-1)$;
at $W$ this uses the nonzero pairing above.
We construct the required subline of $G$. On an ordinary fibre,
constant lattice cocycles identify
\[
 H^1(F_t,\mathcal O_{F_t})
 =\Lambda_{\mathbb C}^*/\operatorname{rowspan}_{\mathbb C}\Pi_c.
\]
Indeed a coboundary among constant cocycles is given by an entire
function with constant period differences. Its derivatives are
periodic holomorphic functions on the compact torus, hence constant;
the function is affine linear. This gives exactly the displayed
kernel, and both sides then have dimension two. The class of
$\gamma=(1,0,0,0)$ is nonzero, since evaluation on columns
$\widehat w,\widehat\delta$ would force both coefficients of a
putative period-row combination to vanish.
The matrices satisfy $\gamma A_j=\gamma$. Translations act trivially
on $H^1(\mathcal O)$, so this holomorphic class extends on both finite
root charts and respects their gluing.

At the cusp $\gamma$ kills
$\Lambda_{\mathrm{tor}}=\mathbb Z\langle\widehat w,\widehat\delta\rangle$.
For local sections of the free toric quotient with changes
$s_i=\Psi_{\lambda_{ij}}s_j$, the constants
$\gamma(\lambda_{ij})$ therefore form a holomorphic Cech cocycle
on the entire neighbourhood. With
$B_0=\left(\begin{smallmatrix}0&1\\-1&0\end{smallmatrix}\right)$,
the edges $e_1,e_2,e_2-e_1$ have deck vectors
$(0,1),(-1,0),(-1,-1)$. The central cocycle values are consequently
$(0,-1,-1)$, which pair to one with $(1,0,-1)$.
Its value on $W$ is nonzero. On the punctured overlap the exponential
map kills precisely $\Lambda_{\mathrm{tor}}$, so its full-lattice
cocycle is exactly $\lambda\mapsto\gamma(\lambda)$, proving equality
with the preceding class. Over the Stein cusp disc, the Leray edge
map $H^1(\bar f^{-1}D_0,\mathcal O_X)\to H^0(D_0,G)$ is an
isomorphism, since $H^1(D_0,\mathcal O)=H^2(D_0,\mathcal O)=0$.
Thus this is a holomorphic, nonvanishing extension of $[\gamma]$.
It supplies a global subline and exact sequence
\[
 0\longrightarrow\mathcal O\xrightarrow{[\gamma]}G
 \longrightarrow\mathcal O(-1)\longrightarrow0.
\]
Again $H^1(\mathcal O(1))=0$ splits it. We have proved
\begin{equation}\label{eq:ab-direct-images}
 R^q\bar f_*\mathcal O_X=
 \begin{cases}
 \mathcal O,&q=0,\\
 \mathcal O\oplus\mathcal O(-1),&q=1,\\
 \mathcal O(-1),&q=2.
 \end{cases}
\end{equation}

\emph{From logarithmic to ordinary tangent cohomology.}
The projection formula gives, for $q=0,1,2$ respectively,
\[
 R^q\bar f_*\mathcal V:
 \quad\mathcal O\oplus\mathcal O(-1),\quad
 \mathcal O\oplus2\mathcal O(-1)\oplus\mathcal O(-2),\quad
 \mathcal O(-1)\oplus\mathcal O(-2).
\]
For $\bar f^*\mathcal O(-5)$ the three expressions are
$\mathcal O(-5)$, $\mathcal O(-5)\oplus\mathcal O(-6)$ and
$\mathcal O(-6)$. Formula \eqref{eq:ab-root-degree} makes every
nontrivial summand in this list acyclic. Leray, whose base has
cohomological dimension one, therefore gives
$h^\bullet(\mathcal V)=(1,1,0,0)$ and zero cohomology for the
horizontal term. Sequence \eqref{eq:ab-log-tangent} yields
$h^\bullet(\mathcal T)=(1,1,0,0)$.

The normalization map $\nu:E\to X$ is an immersion on each local
branch. Let $N_\nu=\operatorname{coker}(T_E\to\nu^*T_X)$.
The exact normal-derivative map is
\begin{equation}\label{eq:ab-normal-quotient}
 0\longrightarrow\mathcal T\longrightarrow T_X
 \longrightarrow\nu_*N_\nu\longrightarrow0.
\end{equation}
For a local equation $z_0\cdots z_{r-1}=0$, it sends
$\sum b_i\partial_{z_i}$ to the $r$ independent normal coefficients
$b_i|_{z_i=0}$. These can be lifted independently, proving
surjectivity; the kernel is exactly the logarithmic tangent module.
This also specifies the global map under changes of branch charts.
The cusp three-form
\[
 \Omega=d\log x_1\wedge d\log x_2\wedge dq
\]
is invariant under $x\mapsto u(q)q^{B_0\lambda}x$ because the
extra logarithmic differentials are multiples of $dq$. In any
unimodular maximal cone the height-one equation is $q=z_0z_1z_2$,
and the determinant-one character change gives
$\Omega=\pm dz_0\wedge dz_1\wedge dz_2$.
It therefore trivializes $K_X$ on the entire cusp neighbourhood,
including every branch intersection. Taking determinants of the
immersion tangent sequence now gives the explicit line identification
$N_\nu=K_E\otimes(\nu^*K_X)^{-1}\simeq K_E$.
Since $E$ is the rational projective surface $\mathrm{dP}_6$,
Serre duality gives $h^\bullet(E,K_E)=(0,0,1)$.
The cohomology sequence of \eqref{eq:ab-normal-quotient} proves
$h^\bullet(X,T_X)=(1,1,1,0)$, as asserted.

\emph{The reduced deformation germ.}
A.4.2 gives a nonzero Kodaira--Spencer vector at this same $c$;
the just computed $h^1=1$ makes its map an isomorphism.
Embed the actual Kuranishi germ in its one-dimensional tangent space
and write its local analytic algebra as $\mathbb C\{t\}/I$.
The existing proper family has a local classifying map, whose pullback
sends $t$ to $g(s)$ with $g'(0)\ne0$. The holomorphic inverse-function
theorem makes $\mathbb C\{t\}\to\mathbb C\{s\}$, $t\mapsto g(s)$,
an isomorphism. Every element of $I$ maps to zero, so $I=0$.
This proves smoothness and reducedness even though $h^2(T_X)=1$.
The classifying map is consequently a local isomorphism of germs;
completeness of the Kuranishi deformation then gives the claimed
miniversality and exhaustion of sufficiently small deformations.
\end{proof}

\paragraph{Scope and a precise changed-input test.}
The hypotheses of A.5.1--A.5.2 fix the original period system and
multiplicities $(3,4)$; B.2 treats its arbitrary affine and cusp data,
and B.3 classifies its stipulated higher-multiplicity variants.
A concrete extension beyond these hypotheses begins with
$y^2=x^3+tx+1$ in \cite[Sections 2--3 and Proposition 6.1]{VTwo26}.
Already one rank-one filling over a specified root $p_1$ of
$4t^3+27=0$ replaces the hexagonal normalization calculation by
\[
 \begin{gathered}
 E'=\mathbb C^*/\langle c_0\rangle,\qquad
 R=\mathbb P_{E'}(\mathcal O_{E'}\oplus L),\\
 R=(\mathbb P^1\times\mathbb C^*)/
       \langle(z,v)\mapsto(\gamma_0z,c_0v)\rangle.
 \end{gathered}
\]
with its boundary sections identified by $v\mapsto\beta_0v$.
Here $c_0,\gamma_0,\beta_0$ are exactly the nonzero central multipliers
in that paper's equation (3.1). The invariant nodal volume form there
trivializes the ambient canonical line, so the branch-normal map has
quotient $\nu_*K_R$ by the same local coefficient map as
\eqref{eq:ab-normal-quotient}. The ruling has
$\pi_*\mathcal O_R=\mathcal O_{E'}$ and $R^1\pi_*\mathcal O_R=0$;
Leray and Serre duality therefore give $h^\bullet(K_R)=(0,1,1)$.
The first unevaluated global datum for transferring the deformation
argument is the connecting map from this single
$H^1(R,K_R)=\mathbb C$ summand into $H^2(T_X(-\log D))$, where
$D$ is the union of the three reduced singular fibres of that model.
Its direct attack is to compute the boundary-translation Cech cocycle
using the full holomorphic unit functions in that equation (3.1),
their central values $(\gamma_0,c_0,\beta_0)$, and the corresponding
vertical frames.
The global identification argument also requires the integral
centralizer and collar defects of that model's four monodromy matrices.
These are specified substitutions to perform for the stated one-fibre
test; no classification or deformation conclusion for that changed
period system is included in A.5.1--A.5.2.

\section{B. Discrete data and the geometric parameter}
\subsection{B.1. A typed fibre dictionary}

\begin{proposition}[B.1.1; Proved]\label{st:B.1.1}
Let $E_\tau=\mathbb C/(\mathbb Z+\tau\mathbb Z)$, $P=[\mu]$, and $L_\mu=\mathcal O_{E_\tau}(P-O)$. The holomorphic map
\[
 [\xi,\eta]\longmapsto[\xi,e^{2\pi i\eta}]
\]
identifies $\mathbb C^2/\Pi_c\Lambda$ with the quotient of $L_\mu^\times$ by the lift
$A(\xi,W)=(\xi+6\mu,e^{2\pi i\beta_c}W)$. Engel's generator, in Definition 2.3 of the attached paper, is its inverse, with translation $-6P$ and multiplier $e^{-2\pi i\beta_c}$. The section $P=[\mu]$ is a generator of the Mordell--Weil group of the stated pointed Weierstrass surface.
\end{proposition}
\begin{proof}
First quotient by columns two, three and four. Exponentiation kills column four; columns three and two then give
\[
 (\xi,W)\sim(\xi+1,W)\sim(\xi+\tau,e^{2\pi i\mu}W).
\]
For the odd theta function with periods $1,\tau$, the quotient $\theta(\xi-\mu)/\theta(\xi)$ is a meromorphic section of this line bundle: its factors are $1,e^{2\pi i\mu}$ and its divisor is $P-O$. Sending its divisor section to this quotient is the asserted line-bundle isomorphism. Quotienting the remaining first column gives exactly $A$. On sufficiently small simply connected charts the logarithm of $W$ is a holomorphic inverse; changes of logarithm are column-four periods. Thus the displayed map is a biholomorphism of the two quotient complex spaces, with its inverse explicitly defined by local logarithms. Choosing $A^{-1}$ changes the generator of the same infinite cyclic group and gives precisely Engel's declared inverse-translation convention.

For clarity the elliptic surface in this dictionary is the displayed cubic, rather than only a curve with the same $j$-value. Let $f_{\rm sc}$ be the branch satisfying
\[
 f_{\rm sc}^4=\frac{12t^3(t-1)}{g_2(\tau)},\qquad
 f_{\rm sc}^6=-\frac{8t^4(t-1)^2}{g_3(\tau)},\qquad
 f_{\rm sc}/t\longrightarrow i\sqrt3/\pi.
\]
The identities $g_2^3/(g_2^3-27g_3^2)=t$ make these equations compatible. Their zero and pole orders give a holomorphic branch on the period cover; the indicated cusp sign fixes it. The map $\xi\mapsto(f_{\rm sc}^2\wp_\tau(\xi),\frac12f_{\rm sc}^3\wp'_\tau(\xi))$ is a pointed isomorphism to
$y^2=x^3-3t^3(t-1)x+2t^4(t-1)^2$. Its factors are $f_{\rm sc}(g_1z)=-f_{\rm sc}/\tau$, $f_{\rm sc}(g_2z)=f_{\rm sc}/\tau$; therefore this isomorphism descends over $B^\circ$.

We first verify that $P=[\mu]$ gives a rational section of the resolved elliptic surface. On $B^\circ$ its cubic coordinates descend by the displayed scale and affine laws. At the two finite centres, $\mu_j$ is not a lattice point, while $f_{\rm sc}$ has orders two and one in the root parameter. Its cubic coordinates are therefore holomorphic and invariant, hence holomorphic in the base coordinate. At the cusp $\mu$ is holomorphic and bounded, and the convergent multiplicative expansion of $\wp$ gives meromorphic cubic coordinates with finite pole orders. Even when its exponential meets one, the pole is controlled by its finite vanishing order; it cannot be identically one because the two finite values of $\mu$ differ. The coordinates are thus meromorphic on $\mathbb P^1$ and rational. The closure of its graph on the proper resolved surface is an irreducible curve; its normalization maps finitely with degree one to the smooth base, so is the base itself. This gives the unique section of that surface extending $P$.

We also verify that the specified section is a Mordell--Weil generator using its full logarithm. A rational section of this cubic pulls back to a holomorphic section of the good elliptic family at both finite covering points: multiply its projective coordinates by a common power to remove poles, and use the nonsingular limiting Weierstrass equation. At the cusp put $v=1/t$, $X=v^2x,Y=v^3y$. Its equation is $Y^2=X^3-3(1-v)X+2(1-v)^2$. A section cannot pass through the node $(1,0)$: substitution of $X=1+O(v)$, $Y=O(v)$ leaves the coefficient $+1$ of $v$ on the left-hand side of the defining equation. At every other central point the multiplicative uniformization is a local isomorphism. Explicitly, at $v=0$ it is
\[
 X=1+\frac{12U}{(U-1)^2},\qquad
 Y=\frac{12\sqrt3 U(U+1)}{(U-1)^3},\qquad U=e^{2\pi i\xi},
\]
with inverse $U=(Y/(X-1)+\sqrt3)/(Y/(X-1)-\sqrt3)$ away from $O$; a local parameter at $O$ is $X/Y$. The convergent lattice expansion of $\wp$ gives the same holomorphic map for small $v$. Hence a section has a holomorphic unit $U(v)$ and a bounded logarithm at the cusp. Lifting the section on the simply connected cover gives a global logarithm $\nu$, with this cusp branch. Its only remaining ambiguity is an integer, since a nonzero multiple of $\tau$ is unbounded.

Its two affine laws have integral defects:
\[
 \nu(g_1z)=-\nu/\tau+\alpha_1+\beta_1\tau(g_1z),\quad
 \nu(g_2z)=\nu/\tau+\alpha_2+\beta_2\tau(g_2z).
\]
Composing them and using cusp invariance gives $\alpha_2=-\beta_1$, $\beta_2=\alpha_1+\beta_1$. Put $n=-\alpha_1-2\beta_1$ and $a=\alpha_1+\beta_1$. Substitution shows that $n\mu+a$ has exactly these defects. Their difference is a global section of the first-coordinate sheaf $\mathcal O_B(-p_0)$ from A.3.1, so is zero. Every rational section is therefore $nP$. The same cusp argument and the distinct fixed values of $\mu$ show that $P$ has infinite order. This proves $\operatorname{MW}(S)=\mathbb ZP$. \end{proof}

\begin{proposition}[B.1.2; Proved]\label{st:B.1.2}
The exact set of nonzero parameters for which every ordinary fibre and both normalized finite root fibres of Engel's period-reference quotient are compact complex two-tori is
\[
 0<|u|<e^{-2\pi M}.
\]
This is the parameter domain for the compact family with both smooth finite root models. The condition on ordinary fibres alone at equality is determined by whether a maximizing point of $D_{\rm part}$ is ordinary; this boundary issue is treated in Part C. For an independent reference, the radius is $|\lambda_0|e^{-2\pi M}$.
\end{proposition}
\begin{proof}
The function $N(\xi,W)=|W|\exp(2\pi\operatorname{Im}\mu\,\operatorname{Im}\xi/\operatorname{Im}\tau)$ is invariant under both line-bundle period identifications. The inverse lift multiplies it by $e^{2\pi D_c}$. More directly the real determinant of the four period columns is a nonzero sign times $(\operatorname{Im}\tau)D_c$. Thus the four columns generate a full lattice exactly when $D_c\ne0$, and then the quotient is a compact torus by a closed fundamental parallelepiped. At $D_c=0$ their real span has dimension three; a rank-four lattice in $\mathbb R^4$ cannot result, so the quotient cannot be a compact complex two-torus. The function $D_{\rm part}$ descends continuously to the connected base with the cusp removed and tends to $-\infty$ at the cusp. Its image is therefore $(-\infty,M]$: the maximum is attained, and the intermediate value theorem fills all smaller values. Consequently $D_c$ is nonzero at every base point exactly when $\operatorname{Im}c<-M$. Since $|u|=e^{2\pi\operatorname{Im}c}$, the assertion follows, with the indicated reference scaling.
\end{proof}

\begin{proposition}[B.1.3; Proved]\label{st:B.1.3}
For the original affine logarithm $\mu$, we have $\wp_\tau(\mu)=0$ identically. Equivalently the generator $P=[\mu]$ of B.1.1 is one of the two sections $P_*=(0,\sqrt2t^2(t-1))$ and $-P_*$. This assertion does not use a numerical evaluation of $\mu$ or of $\beta$.
\end{proposition}
\begin{proof}
We give the intersection calculation that identifies the generator. At $t=1$, put $u=t-1$ and blow up the origin. In the $u$-chart $x=uX,y=uY$, the strict equation is
$Y^2=uX^3-3(1+u)^3X+2(1+u)^4$. Its exceptional curve is smooth, and $P_*$ meets it at $(X,Y)=(0,\sqrt2)$. In the other chart $u=xU,y=xY$, the equation is $Y^2=x-3(1+xU)^3U+2(1+xU)^4U^2$. The implicit function theorem solves for $x$ at the common point; the two fibre branches $U=0,x=0$ meet with intersection multiplicity two. These charts cover the tangent conic and resolve the singularity. The two components have multiplicity one; the one met by $P_*$ is the nonidentity component $D$, with $D^2=-2$ from intersection with the full fibre.

At $t=0$, two blow-ups in the successive $t$-charts give
\[
 Y^2=tX^3+3t^2(1-t)X+2t^2(1-t)^2,
\quad
 Y_2^2=t^2X_2^3+3t(1-t)X_2+2(1-t)^2.
\]
The second equation has two smooth multiplicity-one components $E_{2,\pm}$ at $Y_2=\pm\sqrt2$, and $P_*$ meets one of them. The remaining singular point, in the second $X$-chart, has
\[
 V^2=H(X,T),\quad H=X^2T+(2+3X)T^2-(4X+3X^2)T^3+2X^2T^4,
 \qquad t=XT.
\]
Let $T_0(X)$ solve $H_T(X,T_0)=0$. It has $T_0=-X^2/4+O(X^3)$. Taylor's formula gives
$H-H(X,T_0)=(T-T_0)^2h_0$, $h_0(0,0)=2$, and $H(X,T_0)=X^4e(X)$, $e(0)=-1/8$. These identities follow by substitution in the displayed polynomial, not by discarding its higher terms. With $Z=(T-T_0)\sqrt{h_0}$, $a=V-Z$, $b=(V+Z)/e(X)$, the equation becomes exactly $ab=X^4$ through a local analytic isomorphism. Its resolution has charts
\[
 a=U_i^{i+1}V_i^i,\quad b=U_i^{3-i}V_i^{4-i},\quad X=U_iV_i
 \quad(0\le i\le3),
\]
with changes $U_{i+1}=1/V_i,V_{i+1}=U_iV_i^2$. They give a chain of three $(-2)$-curves. These are the charts of the two successive point blow-ups resolving $ab=X^4$, so the resolution is proper. Retaining $t=XT$ in the comparison gives multiplicities $2,3,2$ on that chain. The two new multiplicity-one components meet its ends; the first exceptional curve meets its middle component, and the original strict fibre meets the other end of that first exceptional curve. All these statements follow by substituting the displayed monomials and their orders; the only fibre components of multiplicity one are the strict component and $E_{2,\pm}$.

Order the six nonidentity components by the graph
$\Theta_1-\Theta_2-\Theta_3-\Theta_4$ and $\Theta_2-\Theta_5-\Theta_6$, with $E_{2,\pm}=\Theta_4,\Theta_6$. Their intersection matrix is $-C$, where
\[
 C=\begin{pmatrix}2&-1&0&0&0&0\\-1&2&-1&0&-1&0\\0&-1&2&-1&0&0\\0&0&-1&2&0&0\\0&-1&0&0&2&-1\\0&0&0&0&-1&2\end{pmatrix},\quad
 C^{-1}=\frac13\begin{pmatrix}6&9&6&3&6&3\\9&18&12&6&12&6\\6&12&10&5&8&4\\3&6&5&4&4&2\\6&12&8&4&10&5\\3&6&4&2&5&4\end{pmatrix}.
\]
Multiplication verifies this inverse. Self-intersections follow as well from the complete multiplicity vector and $F\cdot\Theta_i=0$.

The form $dx\wedge dt/y$ is holomorphic and nonvanishing on the finite resolved charts. In the last charts its expression is a unit times $dX\wedge da/a=-dU_i\wedge dV_i$; the earlier blow-up equations cancel the apparent denominators at $y=0$. At infinity it is $-v^{-1}dX\wedge dv/Y$ and has one fibre pole. Thus $K_S=-F$. Adjunction gives $Q^2=-1$ for every section $Q$, since $F\cdot Q=1$. A section meets a unique multiplicity-one component transversely at a smooth fibre point: the pullback of the full base equation to the section has order one, excluding all other possibilities.

Let $v_Q$ be the intersection vector with $\Theta_i$ and let $d_Q=Q\cdot D$. The possibilities are $v_Q=0,e_4,e_6$ and $d_Q=0,1$. Define the rational divisor class
\[
 \varphi(Q)=Q-O-(Q\cdot O+1)F+\sum\Theta_i(C^{-1}v_Q)_i+\tfrac12d_QD.
\]
It is orthogonal to $O,F,\Theta_i,D$ by direct multiplication. The chord-and-tangent rational function on the generic cubic has divisor $(Q)+(R)-(Q+R)-(O)$; extending it to the surface adds only vertical components. Applying the displayed orthogonal projection therefore proves $\varphi(Q+R)=\varphi(Q)+\varphi(R)$. Expansion gives
\[
 h(Q)=-\varphi(Q)^2=2+2Q\cdot O-v_Q^tC^{-1}v_Q-\tfrac12d_Q.
\]
For $Q\ne O$, the nonnegative intersection $Q\cdot O$ gives $h(Q)\ge2-4/3-1/2=1/6$. The visible $P_*$ is disjoint from $O$ everywhere, and the explicit blow-up charts place it on both nonidentity components. Hence $h(P_*)=1/6$. The logarithm argument in B.1.1 already proved $P_*=nP$ for an integer $n\ne0$, without identifying the visible section. Additivity gives $1/6=n^2h(P)\ge n^2/6$, so $n=\pm1$. The pointed cubic map in B.1.1 has $x(P)=f_{\rm sc}^2\wp_\tau(\mu)$ with $f_{\rm sc}\ne0$ on the ordinary locus. This proves the asserted zero identity.
\end{proof}

\begin{proposition}[B.1.4; Proved]\label{st:B.1.4}
The line-bundle comparison of the period family with Engel's family over $B^\circ$ transports its inverse period lift to a global section of $\operatorname{Hom}(t_{6P}^*M_S,M_S)$ with divisor $F_\infty$. Take this transported section for $c=0$ as the reference linearization. Then $\lambda_0=1$ exactly and $u=e^{-2\pi ic}=w^2$. For another independently fixed reference the precise constant is its scalar ratio with this section. With the cusp branch fixed in Part C and $f_{\rm sc}/t\to i\sqrt3/\pi$, the section is $P=(0,-\sqrt2t^2(t-1))$.
\end{proposition}
\begin{proof}
The zero identity $\wp_\tau(\mu)=0$, proved without numerical input in B.1.3, identifies $P$ with the signed visible section. With the cusp branch in Part C and the specified scale, substitution of $U=-(5+2\sqrt6)$ in the nodal uniformization of B.1.1 gives $Y=-\sqrt2$, so $P=(0,-\sqrt2t^2(t-1))$.

To compare the line bundles as families, use their common divisor $P-O$ and rigidify at the zero section. Put $F_\theta(\xi)=\theta(\xi-\mu)/\theta(\xi)$, $r_\theta=\theta(-\mu)/\theta'(0)$, and let $s_{P-O}$ be the canonical meromorphic divisor section. The map
\[
 [\xi,W]\longmapsto\frac{2r_\theta}{f_{\rm sc}}
        W\frac{s_{P-O}(\xi)}{F_\theta(\xi)}
\]
is a holomorphic line-bundle isomorphism, including at $P,O$ after cancellation of the equal divisors. At $O$ it sends the chosen frame to $dx/y$, since $F_\theta\sim r_\theta/\xi$ and $x/y\sim-\xi/f_{\rm sc}$ with the corresponding canonical divisor frame. Two such maps that agree at $O$ differ by a holomorphic scalar on each compact elliptic fibre and hence agree everywhere. This proves descent under the actual period identifications.

Engel's \cite[Lemma 2.1]{E26}, which we have read together with its proof, identifies $\operatorname{Hom}(t_{6P}^*M_S,M_S)$ with $\pi_S^*\mathcal O_B(1)$. The lemma applies because the preceding pointed cubic map and Mordell--Weil calculation identify his $(S,O,P)$ with ours, including the chosen sign. To determine which section is transported, it suffices to compute its order at one generic smooth point of each special fibre, since a section of this pulled-back line is a base section. At a finite centre write $6\mu=p+r\tau+\delta$ with $(p,r)=(4,-2),(3,-3)$ and $\delta/f_{\rm sc}$ holomorphic; this divisibility follows by linearizing the affine laws, whose homogeneous orders are two and one. Evaluate at $\xi=f_{\rm sc}z$ for generic fixed $z$. The factor $(2r_\theta/f_{\rm sc})/F_\theta(f_{\rm sc}z)$ tends to $2z$, a unit. Restoring the lattice shift $p+r\tau$ at the target changes the fibre multiplier to $e^{-2\pi i(\beta_c-r\mu)}$, whose finite limits are $e^{-2\pi i\kappa_1}$ and $e^{-2\pi i\kappa_2}$. The source and target theta ratios remain units. Therefore the base section has order zero at both finite points. At the cusp $f_{\rm sc}$ has a simple pole and $r_\theta$ has a nonzero limit; the two factors $2r_\theta/f_{\rm sc}$ cancel between source and target. The remaining multiplier is
$e^{-2\pi i\beta_c}=q e^{-2\pi i(\beta_c+\tau)}$, of exact order one. Thus the transported section has divisor $F_\infty$. Engel's \cite[Corollary 2.2]{E26} now identifies it with a nonzero scalar multiple of his reference, since two base sections with the same divisor have constant ratio. Changing $\beta$ by $c$ multiplies it exactly by $e^{-2\pi ic}$. With our stated period-reference normalization, this scalar is one. For any independently chosen reference $\phi_{\rm ref}$, the scalar is instead the precisely defined ratio $\lambda_0=\phi_{\rm part}/\phi_{\rm ref}$, evaluated on any ordinary fibre. Scaling the reference by $a\in\mathbb C^*$ changes $\lambda_0$ to $\lambda_0/a$; no scalar value is invariant under that declared choice.
\end{proof}

\subsection{B.2. All affine vectors and cusp sections}

\begin{proposition}[B.2.1; Proved]\label{st:B.2.1}
Let
\[
 w_1=r_1\varepsilon+k_1d,\qquad w_2=r_2\varepsilon'+k_2d,
 \qquad 3\nmid r_1,\quad 2\nmid r_2,
\]
and use an arbitrary holomorphic cusp collar section $\sigma$ in the construction of \cite[Section 6.4]{AC26}. Write $l_0=\gamma(c(\sigma))$ in its positive cusp-winding convention, and set $p=12l_0-4r_1-3r_2$. Let primes denote a second such datum with the same period solution and cusp fan. Then
\[
 X(c;w_1,w_2,\sigma)\simeq X(c';w'_1,w'_2,\sigma')
\]
holds exactly for the following arithmetic conditions:
\[
 n=c'-c\in\mathbb Z,\qquad
 p'=a p\text{ for some }a\in\{1,-1\},\qquad
 n\equiv k_2-k'_2\pmod2.
\]
All maps are over $B$. A map with the indicated $a$ has fibre-linear part $aI_2$. In particular the complete invariants are
\[
 \bigl(|p|,\ e^{-\pi i(c+k_2)}\bigr),\qquad
 |p|\ge1,\quad\gcd(|p|,6)=1.
\]
Every value of this positive integer occurs. The full collar section enters the construction, but its additional analytic part and its second integral winding contribute no further invariant.
\end{proposition}
\begin{proof}
The determinant calculation in A.1.1 was independent of translations, so the base is still intrinsic. The centralizer proof gives $c'=c+n$ and $aI_2$ exactly as before. Put $S_n=I+nE$ and compare all target periods in source coordinates. On the finite root cover the exact deck equation is
\[
 B_j(g_jz)-R_jB_j(z)=\Pi_c(g_jz)\left(\frac{d_j}{m_j}+\ell_j\right),
 \quad d_j=S_nw'_j-a w_j,
 \quad N_j\ell_j=-d_j.                                      \tag{B2.1}
\]
Necessity follows by the normalized-base-change lift already proved. Conversely the norm equation has a rational solution $r_j^*$ of
$(I-A_j)r_j^*=d_j/m_j+\ell_j$, since $\ker N_j$ over $\mathbb Q$ equals the image of $I-A_j$. We can choose $\gamma(r_j^*)=0$: the columns $\widehat u,\widehat w$ already span that image over $\mathbb Q$. Then $B_j=\Pi_cr_j^*$ is holomorphic on the whole root disc and solves the exact equation. Thus the norm condition is sufficient for a local isomorphism, not only for conjugacy on the central real torus.

The computed norm images translate (B2.1) into
\[
 r'_1\equiv ar_1\pmod3,\quad r'_2\equiv ar_2\pmod4,
 \quad ak_2-k'_2-nr'_2\equiv0\pmod2.                       \tag{B2.2}
\]
There is no condition on $k_1$. Since $r'_2$ is odd, the last condition is the displayed parity. On the finite collar the middle translation is
$A_j=B_j+(\log s_j)/(2\pi i)\,\Pi_cd_j$ and has integral defect $\ell_j$. Its $\gamma$ defect is
$\alpha_j=(ar_j-r'_j)/m_j$.

We prove the precise cusp necessity used next. A holomorphic isomorphism between cusp neighbourhoods with fibre-linear part $aI_2$ has, after lifting near a generic smooth central point to the toric covers, the form
\[
 (x,t_c)\longmapsto(t_c^q u(t_c)x^a,t_c),\qquad q\in\mathbb Z^2,
 \quad u(0)\in(\mathbb C^*)^2.                           \tag{B2.3}
\]
Here the assertion concerns this local lift near a generic central point. To prove it, take a source chart near the generic point of its component and a target chart near its image. The etale toric quotient permits a local lift over a simply connected neighbourhood. The torus characters in a chart at the height-one ray $v$ are $x_i=t_c^{v_i}X_i$, with $X_i$ invertible. On $t_c\ne0$ the affine fibre equation gives $x'_i/x_i^a$ independent of the $X$ coordinates. Its expression in the two charts is therefore $t_c^{v'_i-av_i}$ times a holomorphic unit of $t_c$, including zero. This proves (B2.3). Taking logarithms shows that its winding is in $\Lambda_{\rm tor}=\mathbb Z\widehat w+\mathbb Zd$. Changing a lift adds $(I-A_0)\Lambda=\Lambda_{\rm tor}$, so the vanishing of its class in $\bar\Lambda=\Lambda/\Lambda_{\rm tor}$ is intrinsic. This proof supplies the necessary extension test; no assertion about arbitrary punctured translations is used.

We retain the actual collar convention. The map from the cusp piece to the middle family is $T_{-\sigma}G_0^{-1}$. Consequently the middle translation induced by (B2.3) is
\[
 A_0=sq+a\sigma-\sigma'+\eta_0,
\]
with $\eta_0$ holomorphic on the cusp disc after a logarithm of the unit. For the clockwise generator $s\mapsto s-1$, the class of its defect in $\bar\Lambda$ is $-a c(\sigma)+c(\sigma')$. On the other hand every middle integral cocycle satisfies
$\ell_0=-A_0(\ell_1+A_1\ell_2)$, whose $\gamma$ coordinate is $-\alpha_1-\alpha_2$. Equality of these coordinates is exactly
\[
 l'_0-a l_0=-(ar_1-r'_1)/3-(ar_2-r'_2)/4,
\]
or $p'=ap$. This proves the remaining necessity. Notice that $p\equiv-r_1\pmod3$ and $p\equiv r_2\pmod2$, so $p$ is nonzero and coprime to six.

For sufficiency, $p'=ap$ implies the first two congruences of (B2.2) by reduction modulo three and four. The parity supplies the remaining one. Choose integral solutions $\ell_1,\ell_2$ of (B2.1). Their $\gamma$ coordinates already have the required values. Adding any element of $\ker N_1=(I-A_1)\Lambda$ preserves (B2.1). In particular adding an arbitrary integer multiple of $e_1=(0,1,-2,0)^t$ changes the $\widehat u$ coordinate of $\ell_0$ by an arbitrary integer while leaving its $\gamma$ coordinate fixed. We can therefore arrange
$[\ell_0]=-a c(\sigma)+c(\sigma')$ in the entire rank-two group $\bar\Lambda$. Choose lifts of $\sigma,\sigma'$ on the cusp half-plane. Their remaining defect differs from $\ell_0$ by an actual toric lattice vector; an integral $q$ makes $sq+a\sigma-\sigma'$ have exactly defect $\ell_0$.

The finite translations and this cusp reference now define a locally nonempty affine torsor with linear difference sheaf precisely $\mathcal V$ of A.3.1. The terms $a\sigma-\sigma'$ are retained, even when their chosen lifts have essential singularities at the puncture; they belong to the specified collar reference, while differences of allowed extensions are holomorphic $\eta_0$. The proof of $H^1(B,\mathcal V)=0$ gives a global middle translation. The finite maps descend by (B2.1), and the cusp map is $x\mapsto t_c^qu(t_c)x^a$. The fan is preserved by all integer translations and by $y\mapsto-y$; exponentials of the integer period shift agree. Hence this cusp map and its inverse extend over every stratum. The finite overlap is the equality defining $A_j$, and the cusp overlap is the equality defining $A_0$. Thus the maps glue globally.

Finally parity is equivalent to equality of the two displayed exponentials. Passing from $p$ to $-p$ is effected by $a=-1$ and the same proof; for fixed sign the congruences have just been shown sufficient. Every positive integer coprime to six is obtained by choosing $r_1\in\{1,2\}$ and odd $r_2$ in the appropriate residue classes modulo three and four, and then taking $l_0=(p+4r_1+3r_2)/12$. This also proves that the extra collar data have been accounted for, rather than omitted from the statement.
\end{proof}

\begin{proposition}[B.2.2; Proved]\label{st:B.2.2}
For multiplicities $(3,4)$ and $|p|=1$, every admissible discrete member is biholomorphic to one member of A.5.1, uniquely modulo $2\mathbb Z$. At fixed nonzero admissible Engel parameter $u$ there are exactly two isomorphism classes, exchanged by a positive circuit of $u$ about zero. In the convention that Engel's root generator is the inverse of our clockwise one and his circle class is $d$, the identity of the middle quotient sends
\[
 X_{c+1}\longleftrightarrow X(u;a_0-d/3,\ a_1+d/4).
\]
Reversing both declared root generators reverses these two displayed signs.
\end{proposition}
\begin{proof}
Take the canonical target data $(r'_1,r'_2,k'_1,k'_2,l'_0)=(1,-1,0,0,0)$, for which $p'=-1$. Choose $a$ so that $p'=ap$, and an integer $n\equiv k_2\pmod2$. B.2.1 supplies the isomorphism to $X_{c+n}$. The uniqueness is A.5.1. At fixed $u=e^{-2\pi ic}$, its logarithms differ by integers; the invariant $e^{-\pi i(c+k_2)}$ takes exactly the two square roots of $u$. Both occur, already as $X_c,X_{c+1}$, and A.2.1 distinguishes them.

We spell out the integral comparison to Engel's common basis $B_0=(e_1,e_2,e_3,e_4)$, retaining his $C f_4=e_1+e_4$:
\[
 P_{\rm E}=\begin{pmatrix}0&-2&-1&-2\\2&-1&0&-1\\0&2&1&3\\-1&0&0&0\end{pmatrix},\qquad
 T_0P_{\rm E}=P_{\rm E}A_1,\quad T_1P_{\rm E}=P_{\rm E}A_2.
\]
Its determinant is $-1$, $P_{\rm E}d=\delta_0$, $\psi P_{\rm E}=-\gamma$, $P_{\rm E}\varepsilon=-e_4$ and $P_{\rm E}\varepsilon'=\delta_0-Cf_4$. These are exact lattice homomorphism identities, verified independently through Viaclovsky's intermediate basis in \texttt{ab\_geometry.py}. They also agree with the complex periods: at $z_1$ the complex-linear map
$(\zeta_1,\zeta_2)\mapsto(-\zeta_1/(\rho+1),\zeta_2+\zeta_1/3)$ carries $\Pi_cP_{\rm E}^{-1}$ to Engel's columns $(1,0),(\rho-1,0),((1+\rho)/3,1/3),(0,-\kappa_1)$. At $z_2$ his $Cf_4$ has preimage $d-\varepsilon'$ and fixed period $1-\kappa_2$. Thus the comparison fixes the positive circle class and the actual local complex structures, in addition to their monodromy matrices.

For the sign assertion, inverting the concrete clockwise affine generator changes its invariant translation $\Pi_cv_j/m_j$ to $-\Pi_cv_j/m_j$. Thus our dictionary is $a_{j-1}=-P_{\rm E}v_j/m_j$ in Engel's lattice marking (and its period image in the torus). Integer shift gives $\Pi_{c+1}v_j=\Pi_c(v_j+\gamma(v_j)d)$, with $\gamma(v_1)=1,\gamma(v_2)=-1$, proving the two signs. A positive $u$ circuit changes $\log u$ by $2\pi i$ and hence changes $c$ by $-1$, exchanging the square roots. This gives the claimed permutation with its orientation.
\end{proof}

\begin{proposition}[B.2.3; Refuted, with corrected family]\label{st:B.2.3}
The assertion in \cite[Remark 8.1]{E26} of a global proper family over the original punctured $u$-disc for fixed sphere logarithmic data requires a double-cover qualification. The retained family with such data does not descend through $u=w^2$. It does give a proper holomorphic family on each simply connected $u$-open set, and a global proper family over the punctured $w$-disc as in A.3.1.
\end{proposition}
\begin{proof}
Descent through $u=w^2$ would furnish an isomorphism between the two pulled-back fibres at $w$ and $-w$, namely $X_c$ and $X_{c+1}$. A.2.1 excludes that isomorphism by the nonzero class of $d$ in $\Lambda^{A_2}/N_2\Lambda$. This is an exact order-two obstruction on the normalized finite cover. It applies to all sphere data by B.2.1 because $r_2$ is odd.

There is also a direct analytic witness near zero. The relative Albanese invariant of the multiplicity-four reduced fibre in a holomorphic descended family would be a single-valued holomorphic function of $u$. After pullback it would be even in $w$. But A.4.1 gives $j_2(w)=C^{-1}w^{-1}+744+O(w)$, where $C=e^{-\pi i\kappa_2^{\rm part}}\ne0$; its Laurent coefficient of $w^{-1}$ is nonzero. Thus it is not even on any punctured neighbourhood. This verifies the obstruction independently of a choice of isomorphism lift. Local inverse logarithms of $u$ give the proper local families already constructed; coherent even shifts give their global double-cover family. The untwisted middle quotient itself has the integer shear identification, so this refutation concerns the fixed sphere logarithmic data in the stated remark.
\end{proof}

\subsection{B.3. Higher multiplicities}

\begin{proposition}[B.3.1; Proved]\label{st:B.3.1}
Let $k,l$ be positive integers with $\gcd(3k,4l)=1$, and put $m_1=3k,m_2=4l$. Consider Viaclovsky's sphere members expressed in the same clockwise marking, with $w_j=r_j\varepsilon_j+s_jd$, $\gcd(r_j,m_j)=1$, cusp winding $l_0$, and $p=m_1m_2l_0-m_2r_1-m_1r_2=\pm1$, where $\varepsilon_1=\varepsilon,\varepsilon_2=\varepsilon'$. Two such members are biholomorphic precisely when their multiplicity pair agrees and, for some $a\in\{1,-1\}$ and integer $n=c'-c$,
\[
 p'=ap,\quad r'_j\equiv ar_j\pmod{m_j},\quad
 nr'_1\equiv as_1-s'_1\pmod{k},\quad
 nr'_2\equiv as_2-s'_2\pmod{2l}.
\]
The congruences on $r'_j$ already follow from $p'=ap$. For fixed discrete vectors and cusp collar the identification subgroup is $2kl\mathbb Z$. For fixed admissible $u$ and this multiplicity pair, all the stated sphere translation data give exactly $2kl$ isomorphism classes, permuted in one cycle by a circuit of $u$ about zero.
\end{proposition}
\begin{proof}
We first recover the base for these sphere members without using the degree-two canonical pencil. A biholomorphism $F:X\to X'$ gives $g=f'F$. Since $H^2(S^6,\mathbb R)=0$, the closed form $g^*\omega_{\rm FS}$ is exact. Restrict to a regular torus $T$ and wedge with a flat Kahler form on $T$. Stokes' theorem gives integral zero, while the integrand is nonnegative and strictly positive wherever $dg|_T\ne0$. Thus $g$ is constant on each regular fibre. The image of $(f,g)$ is a compact analytic curve whose projection to $B$ has degree one; its normalization gives a holomorphic $h:B\to B$ with $g=hf$. Apply the same construction to $F^{-1}$ to see that $h$ is an automorphism. The three homological local monodromy orders are three, four and infinity, independently of $k,l$; their intertwining maps under $F$ therefore fix each special point. Hence $h=1$. The nonreduced fibre multiplicities now force equal $k,l$. This is the full reduction in the read Proposition 4.7 of Viaclovsky's paper, written here with its exact maps.

The root model of degree $m_1$ is the pullback of the degree-three model by $s\mapsto s^k$, and likewise at the other point. Its linear action remains $A_1$ or $A_2$, while its root generator has order $m_j$. The primitive invariant coordinate $r_j/m_j$ proves freeness of every nonidentity power. The normalization and deck-commutation arguments of A.1.2 and A.2.1 apply to these root covers without changing a generator. Their norms are
\[
 N_{3k}=kN_1,\qquad N_{4l}=lN_2,
\]
so the fixed-coordinate norm images are $3k\mathbb Z\varepsilon+k\mathbb Zd$ and $4l\mathbb Z\varepsilon'+2l\mathbb Zd$. The norm equation (B2.1), with $m_j$ replaced by these actual covering degrees, gives exactly the displayed congruences. The centralizer and first-column comparison remain unchanged.

For completeness the integer $p$ records the primitive sphere relation in this marking. The common torus coinvariants are detected by $\gamma$, since $(A_1-I)\Lambda+(A_2-I)\Lambda=\ker\gamma$ by their explicit saturated bases. The finite meridian relations are $x^{m_1}=c^{r_1}$, $y^{m_2}=c^{r_2}$, and the cusp relation is $xy=c^{l_0}$ with $c$ central. Eliminating $y$ gives an abelian relation matrix with determinant $-p$ and primitive first elementary divisor because $\gcd(m_1,m_2)=1$. Thus the sphere data have $p=\pm1$. Reducing $p'=ap$ modulo $m_1,m_2$ gives the two congruences for $r_j$.

To prove sufficiency, use the local rational solutions of the norm equation and the identical cusp comparison in B.2.1. The first-coordinate extension on the new root discs has vanishing orders $2k$ and $l$, precisely the orders of the pulled-back $\Theta$. Dividing by it again gives $\mathcal O_B(-p_0)$, and finite averaging again gives the same rank-two extension sheaf. Thus $H^1(B,\mathcal V)=0$ still removes the analytic cocycle. The images of $I-A_j$ and their arbitrary $\widehat u$ adjustments are unchanged, so $p'=ap$ is exactly the remaining cusp winding condition. Every local map and overlap is the one already written in B.2.1, with the new denominator $m_j$; this proves global sufficiency.

For identical data, $p\ne0$ excludes $a=-1$. The congruences reduce to $k\mid nr_1$ and $2l\mid nr_2$. The two primitivity assumptions imply $k\mid n$ and $2l\mid n$. Since $\gcd(k,2l)=1$, their least positive common multiple is $2kl$. For different sphere translation data, $r'_1$ is invertible modulo $k$, $r'_2$ modulo $2l$, and the Chinese remainder theorem gives exactly one integer shift class modulo $2kl$. The values of $s_1$ modulo $k$ and $s_2$ modulo $2l$ all occur without affecting freeness or $p$. Thus there are exactly $2kl$ classes at fixed $u$. The change $c\mapsto c-1$ on a positive $u$ circuit acts by a generator of this cyclic residue group. Program \texttt{ab\_exact.py} verifies the repeated norms and shift steps for five coprime pairs; the preceding symbolic identities prove all pairs.
\end{proof}

\subsection{B.4. Automorphisms and the remaining real-structure calculation}

\begin{proposition}[B.4.1; Proved]\label{st:B.4.1}
For each member of the retained $(3,4)$ family, $\operatorname{Aut}(X_c)$, including all its components, is $\mathbb C/\mathbb Z\simeq\mathbb C^*$. Its action on the ordinary fibres is $\zeta\mapsto\zeta+be$, and the exponential group isomorphism is $b\mapsto e^{2\pi ib}$.
\end{proposition}
\begin{proof}
A.1.1 and A.2.1 give base identity, $n=0$ and fibre-linear part $I$. Let $\ell_1,\ell_2$ be the integral defects of a lift of its translating section to the universal cover of $B^\circ$. Finite extension gives $N_j\ell_j=0$. The saturated image calculations in A.1.2 give the integral equality $\ker N_j=(I-A_j)\Lambda$. The cusp valuation argument (B2.3) gives $\ell_0\in\Lambda_{\rm tor}$. Write
\[
 \ell_1=ae_1+bf_1,\qquad \ell_2=ce_2+df_2.
\]
The $\widehat u$ coordinate of $\ell_1+A_1\ell_2$ is $a-c-2d$, so the cusp condition is exactly $a=c+2d$. Changing the global lift by a period $\lambda$ changes this pair by the simultaneous coboundary $((I-A_1)\lambda,(I-A_2)\lambda)$. In the four displayed coefficient coordinates its matrix is
\[
 C=\begin{pmatrix}-6&1&-1&0\\2&-1&0&0\\0&1&1&0\\-3&0&-1&0\end{pmatrix}.
\]
The lattice of all allowed pairs has basis matrix
$K=\left(\begin{smallmatrix}0&1&2\\1&0&0\\0&1&0\\0&0&1\end{smallmatrix}\right)$. The first three columns of $C$ expressed in this basis have determinant one. Hence every allowed pair is an integral coboundary; the computation is exact, including its saturation, and is independently checked by Hermite form in \texttt{ab\_exact.py}. After changing the lift, all its defects vanish.

The corrected lift extends as an equivariant holomorphic vector on the finite discs. A holomorphic local lift exists on the whole root disc by lifting the translating section through its covering vector bundle. Its defect is in $\ker N_j=(I-A_j)\Lambda$; subtracting an integral period kills that defect while preserving holomorphy. The corrected global lift and this local defect-zero lift differ by an invariant period, which is holomorphic through the centre. At the cusp, the local toric logarithm in (B2.3) is $sq+\eta(t_c)$. The chosen global lift differs from it by $\Pi_c(s)\lambda=sB_0\bar\lambda+C(t_c)\bar\lambda+\lambda_{\rm tor}$. Defect zero gives exactly $q+B_0\bar\lambda=0$, so the logarithmic terms cancel and the remaining vector is holomorphic. This proves the extension at the cusp as well. We therefore obtain a global section of $\mathcal V$. Its exact sequence in A.3.1 gives $H^0(B,\mathcal V)=\mathbb Ce$. Conversely each such section defines the displayed translations on the finite covers and torus multiplication at the cusp, hence a global automorphism. Its kernel consists of global invariant periods. The common fixed lattice is $\mathbb Zd$, and $\Pi_cd=e$, giving kernel exactly $\mathbb Z$. This proves the full group assertion and its connectedness.
\end{proof}


\begin{proposition}[B.4.2; Proved]\label{st:B.4.2}
Use the normalization $(\beta^{\rm part}+\tau)(p_0)=0$. There is an antiholomorphic diffeomorphism $X_c\to X_{c'}$ exactly when
\[
 c'-(1-\bar c)\in2\mathbb Z,
 \qquad\text{equivalently }w'=-\bar w.
\]
Consequently $\overline{X_c}\simeq X_{1-\bar c}$. The members admitting antiholomorphic self-maps are exactly those with $\operatorname{Re}c\in\tfrac12+\mathbb Z$, equivalently $w=-\bar w$. On each such member the group of holomorphic and antiholomorphic diffeomorphisms is
\[
 \operatorname{Aut}^{\pm}(X_c)\simeq\mathbb C^*\rtimes\mathbb Z/2,
 \qquad z\longmapsto1/\bar z
\]
for the action of the second factor. Its antiholomorphic maps form the nonidentity coset. For all other members the full group is the holomorphic group of B.4.1.
\end{proposition}
\begin{proof}
We first construct the comparison before returning its affine data to the canonical choice. Lift conjugation of the real orbifold base to the antiholomorphic involution $\iota$ of its universal cover whose cusp branch has $\tau(\iota z)=-\overline{\tau(z)}$. A lift exists by the universal covering property; it extends at the finite orbifold points, whose local equations are $s_j^{m_j}=t_j$. Choose the usual paths from a real point $t>1$, so reflection fixes the second centre and sends the first chosen lift to $g_2^{-1}z_1$. The induced involution of the deck group is
\[
 \phi(g_1)=g_2^{-1}g_1^{-1}g_2,\qquad
 \phi(g_2)=g_2^{-1},\qquad \phi(g_0)=g_0^{-1}.
\]
These identities also follow by reflecting the lifted peripheral paths; their product and their squares agree with $g_1g_2g_0=1$. The cusp branch fixes the lift uniquely.

All reflection identities for the periods can be checked before invoking uniqueness. Denote their affine transformation formulas by $G_1,G_2$ on the triple $(\tau,\mu,\beta)$ and put $K(\tau,\mu,\beta)=(-\bar\tau,1-\bar\mu,-\bar\beta)$. Direct rational substitution gives
\[
 KG_1K=G_2^{-1}G_1^{-1}G_2,\qquad KG_2K=G_2^{-1}.
\]
For example $G_2^{-1}(\tau,\mu,\beta)=(-1/\tau,(1-\mu)/\tau,\beta+3-6(1-\mu)^2/\tau)$ and
$G_1^{-1}(\tau,\mu,\beta)=(1/(1-\tau),1-\mu/(1-\tau),\beta-2+6\mu^2/(1-\tau))$. These formulas verify every constant in the two equalities. The reflected $\mu$ is cusp bounded and has the same affine laws, so their difference is a global section of $\mathcal O_B(-p_0)$ and vanishes. After this equality, the difference between the reflected and original $\beta^{\rm part}$ is a holomorphic invariant base function, hence constant; its cusp value is zero under the stated normalization. Therefore
\[
 \mu(\iota z)=1-\bar\mu(z),\qquad
 \beta^{\rm part}(\iota z)=-\overline{\beta^{\rm part}(z)}.
\]
The original $\tau$ identity follows from its $j$ equation and its specified cusp branch. In particular $D_{-\bar c}(\iota z)=D_c(z)$, so this comparison stays in the exact admissible domain.

Let
\[
 H=\begin{pmatrix}1&0&0&0\\0&1&0&0\\-6&0&-1&0\\0&-1&0&-1\end{pmatrix}.
\]
Then $H^2=I$, $HA_2H=A_2^{-1}$, $HA_1H=A_2^{-1}A_1^{-1}A_2$, and
\[
 \Pi_{-\bar c}(\iota z)H=-\overline{\Pi_c(z)}.
\]
Thus $(z,\zeta)\mapsto(\iota z,-\bar\zeta)$ is a well-defined antiholomorphic torus map. It also descends under the middle deck action. Indeed covariance and the displayed lattice conjugacies imply $R_{\phi(g)}(\iota z)=\overline{R_g(z)}$ by multiplying the full real-invertible period matrices. This proves exact compatibility, not merely an equality of fibre lattices.

At the first finite centre compose this map with the target deck identification $g_2$ to return the centre $g_2^{-1}z_1$ to $z_1$. Its lattice map becomes $A_2H$ and it conjugates the clockwise generator to the inverse generator. At the second centre its lattice map is $H$ with the same reversal. Consequently the target affine vectors are exactly
\[
 w'_1=-A_2H\varepsilon=-\varepsilon+d,\qquad
 w'_2=-H(-\varepsilon')=\varepsilon'-3d.
\]
The minus signs arise because the inverse of a generator with invariant translation $w'_j/m_j$ has translation $-w'_j/m_j$. This calculation verifies the full affine actions on the root covers. For the overlaps choose the root coordinates of the real local base coordinates $t$ and $t-1$ compatible with these local reflections. This choice has an explicit comparison with any original root coordinate: their ratio is a holomorphic unit, so changing it changes the collar by $(2\pi i)^{-1}\log h\,\Pi w_j$, a translation extending on the entire root family. Its inverse gives a holomorphic isomorphism of the two glued descriptions. Thus the extension property and final classification transfer to the original coordinates. Then $\log s'_j=\overline{\log s_j}$ on paired branches. Under $\zeta\mapsto-\bar\zeta$, conjugation of $1/(2\pi i)$ supplies the compensating sign, and the identity of the two logarithmic collar translations follows exactly. Two roots of the same fixed real coordinate differ by an $m_j$th root of unity; this changes the comparison only by a constant period-coordinate translation, and the additive correction of A.3.1 glues those constants. The preceding holomorphic-unit comparison handles a different original local base coordinate. Its inverse has the same construction.

At the cusp choose $s=(2\pi i)^{-1}\log t_c$, so $s\mapsto-\bar s$. The full holomorphic unit matrix satisfies
\[
 C_{-\bar c}(\bar t_c)=-\overline{C_c(t_c)}+
       \begin{pmatrix}6&0\\0&1\end{pmatrix}.
\]
The integral extra matrix has trivial exponential. Hence $(x,t_c)\mapsto(\bar x,\bar t_c)$ intertwines the full cusp deck groups, with the same integral index $\lambda$. The toric fan is defined over $\mathbb Z$, so this map extends antiholomorphically on every toric chart and stratum. It preserves the untwisted cusp collar. We have therefore constructed an antiholomorphic map to the actual target datum with parameter $-\bar c$, $r'_1=-1,k'_1=1,r'_2=1,k'_2=-3,l'_0=0$, whose $p'$ is $+1$.

Apply B.2.1 to normalize this target back to the canonical $p=-1$ datum. Take $a=-1$ and the odd shift $n=1$. This gives a holomorphic map to $X_{1-\bar c}$. Every other antiholomorphic map, composed with the inverse of this one, is a holomorphic comparison between canonical members. A.5.1 gives exactly the stated congruence. Exponentiating it gives $w'=-\bar w$.

It remains to verify the group extension rather than infer its square from $H^2$. For a real member, choose any resulting antiholomorphic self-map $J$. Its derivative on ordinary fibre coordinates is $+\text{complex conjugation}$: the initial reflection has derivative $-\text{complex conjugation}$, while the normalization just used has linear part $-I_2$. Hence it conjugates translation by $b\in\mathbb C/\mathbb Z$ to translation by $\bar b$. Its square is a holomorphic translation $d$. Commuting $J$ with $J^2$ shows $d=\bar d$ in $\mathbb C/\mathbb Z$. This equality forces the imaginary part of a representative of $d$ to vanish, so take $d\in\mathbb R$. Replacing $J$ by the composition of translation by $-d/2$ with $J$ changes its square to $d-d/2-d/2=0$. We have an antiholomorphic involution. Under $b\mapsto e^{2\pi ib}$, its conjugation action is $z\mapsto1/\bar z$. This proves the split group and specifies its antiholomorphic coset.
\end{proof}
